\documentclass[
  a4paper, 
  10pt, 
  parskip=half, 
  DIV=12,
  toc=bibliography,
]{scrartcl}

% \usepackage{showframe}

\usepackage{polyglossia}
\setdefaultlanguage{spanish}

% \usepackage{multicol}

% \usepackage{subfiles} 
% \usepackage{tocbibind}
% \usepackage[toc]{appendix}

% \usepackage[justification=centering]{caption}
% \usepackage{subcaption}

% \usepackage[shortlabels]{enumitem}
% \usepackage{epigraph}

\usepackage{fontspec}
% switch to a monospace font supporting more Unicode characters
\setmonofont{JetBrains Mono}[
    Contextuals = Alternate,
    Ligatures = TeX,
    Scale=MatchLowercase
]

\usepackage{minted}
\usemintedstyle{pastie}

% Use KOMA-Script's nice headers instead of fancyhdr
% \usepackage[Conny]{fncychap}

% === Math facilities ===

% Theorem environments
\usepackage{amsthm}

\theoremstyle{definition}

\newtheorem{theorem}{Teorema}[section]
% \newtheorem{proposition}[theorem]{Proposition}
% \newtheorem{corollary}[theorem]{Corollary}
% \newtheorem{lemma}[theorem]{Lemma}
\addto\captionsspanish{\renewcommand\proofname{Demostración}}

\newtheorem{definition}[theorem]{Definición}
\newtheorem{example}[theorem]{Ejemplo}
\newtheorem{problem}[theorem]{Ejercicio}

% \newtheorem{remark}[theorem]{Remark}
\newtheorem{notation}[theorem]{Notación}

\usepackage{centernot}
\usepackage{amssymb}
\usepackage{mathtools}
\usepackage{tikz-cd}

\newcommand{\funfont}[1]{\text{#1}}
\newcommand{\tyfont}[1]{\texttt{#1}}
\renewcommand{\implies}{\rightarrow}
\renewcommand{\impliedby}{\leftarrow}
\renewcommand{\iff}{\leftrightarrow}
\renewcommand{\mapsto}{\Rightarrow}
\newcommand{\Prop}{\mathbf{Prop}}
\newcommand{\placeholder}{\square}

\newcommand{\ImpIfElse}[3]{\texttt{if} \ {#1} \ \{{#2}\} \ \texttt{else} \ \{{#3}\}}
\newcommand{\ImpWhile}[2]{\texttt{while} \ {#1} \ \{{#2}\}}
\newcommand{\ImpConcat}[2]{{#1} \texttt{;} \ {#2}}
\newcommand{\ImpSkip}{\texttt{skip}}
\newcommand{\ImpAssign}[2]{{#1} = {#2}}

\newcommand{\LeanIfElse}[3]{\text{if} \ {#1} \ \text{then} \ {#2} \ \text{else} \ {#3}}

\newcommand{\ImpPlus}{\mathbin{\texttt{+}}}
\newcommand{\ImpMinus}{\mathbin{\texttt{-}}}
\newcommand{\ImpTimes}{\mathbin{\texttt{*}}}

\newcommand{\ImpNot}{\texttt{!}}
\newcommand{\ImpAnd}{\mathbin{\texttt{\&\&}}}
\newcommand{\ImpOr}{\mathbin{\texttt{||}}}
\newcommand{\ImpEq}{\mathbin{\texttt{==}}}
\newcommand{\ImpLe}{\mathbin{\texttt{<=}}}


\usepackage{scalerel}

\DeclareMathOperator*{\Forall}{\scalerel*{\forall}{\sum}}
\DeclareMathOperator*{\BigPi}{\scalerel*{\prod}{\sum}}
\DeclareMathOperator*{\BigSigma}{\scalerel*{\sum}{\sum}}

% === End math facilities ===

\usepackage{csquotes}

\usepackage{bookmark}
\hypersetup{
    colorlinks=true,
    % linkcolor=blue,
    % filecolor=magenta,      
    % urlcolor=cyan,
    % pdftitle={Overleaf Example},
    % pdfpagemode=FullScreen,
    }

% Los apéndices se llaman apéndices
% \renewcommand\appendixpagename{Apéndices}
% \renewcommand{\appendixtocname}{Apéndices}

% mainmatter no reinicia el contador de páginas
% \makeatletter
% \renewcommand\mainmatter{
%     \@mainmattertrue\cleardoublepage\renewcommand\thepage{\arabic{page}}}
% \makeatother

% \renewcommand{\thesubsection}{\thesection.\alph{subsection}}

% \providecommand{\keywords}[1]
% {
%   \small	
%   \textbf{\textit{Palabras clave---}} #1
% }

% define math symbol for big step arrow
\newcommand{\bigstep}{\Longrightarrow}

\title{Presentación con Lean 4 de algunos Resultados clásicos sobre
Semántica de Lenguajes de Programación}
\author{Ricardo Maurizio Paul}
% \date{2020--2024}

\begin{document}

\maketitle

\begin{abstract}
  La utilización de Asistentes para la Demostración de resultados matemáticos parece que terminará por
  constituirse en una absoluta necesidad a fin de poder ratificar la corrección de pruebas previamente
  realizadas “manualmente”. Además, la incorporación de esas pruebas formales a los bancos de la
  aplicación con la que se hayan realizado permitirá tanto la utilización de los resultados probados con total
  garantía de calidad, como el cotejo de los procedimientos utilizados en su obtención, que darán lugar a
  metodologías para el futuro desarrollo de nuevas pruebas a través de patrones o “tácticas”, como se les
  suele llamar en el argot.
  La formalización de las Semánticas de los Lenguajes de Programación constituye un interesante campo de
  aplicación de estos desarrollos por varias razones. En una primera aproximación parece un campo
  particularmente adecuado, pues las citadas semánticas se definen con elementos lógicos y algebraicos
  básicos que auguran una relativamente fácil formalización total. Sin embargo, cuando se entra en el detalle
  aparecen dificultades inesperadas cuya resolución sacará a la luz las técnicas necesarias para tratar con
  éxito casos como estos. Al tiempo, pese a que las componentes matemáticas son evidentes, de nuevo la
  forma concreta en que aparecen y se combinan se separa un tanto de la habitual en campos puramente
  algebraicos, por lo que se han de crear nuevas técnicas para lidiar con ellas.
  Como novedad utilizaremos en este caso la nueva versión 4 del lenguaje Lean, y una parte del trabajo se
  dedicará a presentar las novedades con respecto a la versión 3, con la que se han hecho ya algunos trabajos
  previos en la Facultad. Utilizándola se tratará de llegar lo más lejos posible en la formalización de los
  resultados que se han estudiado en nuestro curso de Teoría de Programación, aunque buscando más la
  explicación detallada de las técnicas utilizadas en las distintas pruebas, que la mera acumulación de
  pruebas. Todas las demostraciones que se hagan se recopilarán en una entrada bien organizada que
  eventualmente se tratará de aportar al repositorio abierto que la organización que soporta Lean facilita
  totalmente en abierto a los usuarios del lenguaje.
\end{abstract}
\pagebreak

\tableofcontents
\pagebreak

\section{Introducción}

Los asistentes de pruebas (o \emph{proof assistants}) son sistemas de software que ayudan en la
construcción y verificación de demostraciones matemáticas mediante el uso de
ordenadores. La creación de asistentes de pruebas se remonta a los años 50 y 60, con
el sistema \emph{Logic Theorist}\footnote{\url{https://en.wikipedia.org/wiki/Logic_Theorist}} (1956). 
Este el primer programa diseñado para realizar \enquote{razonamiento automático} y
fue capaz de demostrar teoremas de lógica proposicional y predicados. 
Su creación precede a la teoría de tipos moderna y 
el isomorfismo de Curry-Howard por lo que usaban una
lógica diferente basada en la \emph{deducción natural} y la \emph{lógica de primer orden}
(\emph{first-order logic} o \emph{FOL}).\@ El sistema fue capaz de demostrar teoremas
interesantes pero tenía un alcance limitado. A grandes rasgos el sistema
generaba un \emph{árbol de búsqueda} de todas las posibles demostraciones y
las evaluaba para encontrar una demostración válida, de forma similar a 
Prolog.

El sistema \emph{Automath} (1967) de De Bruijn fue uno de los primeros
sistemas en usar una teoría de tipos y el isomorfismo de Curry-Howard, que fue 
descubierto independientemente por De Bruijn.

En los años 70 se desarrollaron sistemas más potentes, como \emph{Logic for Computable Functions}, 
basado en la \emph{Lógica de Funciones Computables} o \emph{LCF} de Dana Scott, este fue uno de los
primeros sistemas en usar pruebas por inducción, una característica fundamental
para razonar sobre funciones definidas recursivamente y una de las técnicas 
fundacionales de los asistentes de pruebas modernos. La creación del
lenguaje de programación ML, el padre de todos los lenguajes funcionales, también
se atribuye al desarrollo de LCF, ya que fue diseñado para escribir \emph{tácticas}
en este demostrador. El enfoque moderno de construcción de pruebas \enquote{paso a paso},
usado por todos los asistentes modernos, también se remonta a LCF.\@

En los años 70 también se refinaron ideas como la teoría de tipos, en
particular la teoría de tipos de Martin-Löf, y el isomorfismo de Curry-Howard,
dos conceptos clave que son la base de todos los asistentes de pruebas modernos.
En los años 80 estas ideas fueron implementadas en sistemas como \emph{Nuprl} y
\emph{Coq}. 

Nuprl es uno de los primeros sistemas en usar una teoría de tipos
constructiva en su núcleo, donde las demostraciones son objetos de primera clase,
enfatizando su contenido computacional y permitiendo la \emph{extracción de
código} de programas a partir de demostraciones verificadas. Nuprl también
introdujo tipos dependientes, una característica poderosa que permite que el
tipo de un valor dependa del propio valor. Con este sistema se puede codificar
y demostrar una gran cantidad de la teoría matemática existente.

Un poco más tarde llegó Coq, un sistema basado en el CIC (Cálculo de Construcciones
Inductivas), una teoría de tipos que extiende el Cálculo de Construcciones 
(\emph{Calculus of Constructions} o \emph{CoC}) con tipos inductivos. Coq es uno de los
asistentes de pruebas más populares, y se utiliza ampliamente en la academia y
la industria hasta el día de hoy. Veremos que el CIC también es la base de Lean, un
asistente de pruebas más reciente que ha ganado una popularidad significativa en
los últimos años.

Mientras tanto, se desarrolló otra línea de investigación, la familia de
asistentes de pruebas basados en \emph{HOL} (\emph{Higher Order Logic}), una lógica que
extiende la lógica de primer orden con cuantificadores de orden superior. Estos
sistemas están basados en la arquitectura \emph{LCF} y utilizan un \emph{núcleo} para
asegurar la corrección de las pruebas. El sistema más popular de esta familia es
\emph{Isabelle}, un asistente de pruebas tan potente como Coq pero con una
fundamentación lógica diferente. 

Esto nos lleva a Lean, el asistente de pruebas que usaremos en este texto. Lean
es un asistente de pruebas moderno desarrollado por Leonardo de Moura en Microsoft
Research y basado en el CIC.\@ Podemos ver Lean como una evolución de Coq, con una
sintaxis más moderna, y un sistema de tipos más potente, más adelante discutiremos
las diferencias en más detalle. Otra característica interesante de Lean es que 
su comunidad ha creado una biblioteca de matemáticas muy rica, conocida como \emph{mathlib}, 
donde podemos encontrar resultados sobre teoría de números, álgebra, teoría de
categorías, cálculo, y mucho más. 

Dejando de lado los sistemas basados en HOL, todos estos sistemas están basados en el
isomorfismo de Curry-Howard, un concepto fundamental que relaciona demostraciones
y programas. Para entenderlo primero debemos introducir los conceptos básicos de
la teoría de tipos.

\section{Teoría de tipos}

La teoría de tipos es una rama de la lógica matemática que se ocupa del estudio
de los tipos, una abstracción similar en algunos aspectos a los conjuntos, que 
serán más familiares para el lector. La teoría de tipos surgió a principios del siglo XX
para resolver la crisis de los fundamentos de las matemáticas, una serie de
problemas que surgieron a raíz del descubrimiento de paradojas en la teoría de
conjuntos, como la paradoja de Russell.

De la misma forma que los \emph{conjuntos} son una noción primitiva en la teoría de conjuntos,
los \emph{tipos} junto con las \emph{funciones} son una noción primitiva en la teoría de tipos.
La idea principal detrás de los tipos es clasificar los objetos matemáticos de
forma que las funciones solo puedan aplicarse a objetos de un tipo determinado,
este concepto se conoce como \emph{seguridad de tipos} (\emph{type-safety}). 
Los programadores reconocerán este concepto, ya que es la base de los sistemas de tipos
estáticos en lenguajes de programación como Haskell, Rust y Scala.

Como introducción informal consideramos el siguiente ejemplo: definimos un tipo
\(\mathbb{N}\) que representa los números naturales, y definimos una función
\(\texttt{add}: \mathbb{N} \to \mathbb{N} \to \mathbb{N}\) que toma dos números
naturales y devuelve su suma\footnote{Más adelante formalizaremos la sintaxis de los tipos y funciones.}. 
Esto asegura que la función \(\texttt{add}\) solo puede aplicarse a valores de 
tipo \(\mathbb{N}\), y no a ningún otro tipo y que el resultado de la función
también es de tipo \(\mathbb{N}\).

Históricamente los orígenes de la teoría de tipos se remontan a múltiples
autores como Church, Curry, Russell, y otros. Pero la teoría de tipos moderna se basa en el
trabajo de Pier Martin-Löf, quien desarrolló una teoría de tipos que combina la
teoría de tipos clásica con el \(\lambda\)-cálculo. De esta manera las
\enquote{buenas} propiedades de las teorías de tipos anteriores reciben una
\emph{interpretación computacional}, a través del isomorfismo de Curry-Howard
que nos permite razonar sobre demostraciones matemáticas. 

La teoría de tipos de Martin-Löf es \emph{intuicionista}, es decir, las
demostraciones no son solo un argumentación lógica sino un método para
construir un \emph{testigo} de la verdad de una proposición. Esto contrasta con la lógica
clásica, donde las demostraciones son solo una secuencia de pasos lógicos que
muestran la verdad de una proposición. En la lógica intuicionista, el principio del
tercio excluido no es válido, como tampoco lo es la \emph{reducción al absurdo} (con
algunas salvedades).

\subsection{Lambda cálculo}

Un desarrollo detallado del \(\lambda\)-cálculo se puede encontrar en~\cite{barendregt1984lambda},
pero por completitud haremos una breve introducción.
El \(\lambda\)-cálculo es un sistema formal para expresar \emph{computación} basado en la
abstracción y aplicación de funciones. 
Fue creado por Alonzo Church en la década de 1930 como una forma de estudiar los fundamentos de las matemáticas y
la computación y desde entonces se ha convertido en parte fundamental de la informática.
El cálculo lambda es, en esencia, un \enquote{lenguaje de programación} con solo un tipo de dato, la función, y
dos operaciones, la abstracción y la aplicación.

En el \(\lambda\)-cálculo, las funciones son vistas como \emph{reglas} en vez de
\emph{relaciones}, esto enfatiza el aspecto computacional de las funciones, y
nos permite razonar sobre la computación de una forma más abstracta. Esto contrasta
con la teoría de conjuntos, donde las funciones son vistas como \emph{gráficas}.

La sintaxis del \(\lambda\)-cálculo es simple, consiste en variables (\(x, y, z\)),
abstracción (\(\lambda x \mapsto M\)), y aplicación (\(M \ M\)), en 
\emph{forma de Backus-Naur} (\emph{BNF}):
\begin{align*}
  M \Coloneqq x & \mid \lambda x \mapsto M \mid M \ M
\end{align*}
La notación (\(\mapsto\)), en vez del tradicional (.) no es accidental, y revela una conexión
profunda entre la abstracción de funciones y la implicación lógica, una conexión que exploraremos más adelante.

\begin{example}\label{ex:lambda_expr}
Podemos aplicar la función \emph{sucesor} al número \(3\) como:
\(
(\lambda x \mapsto x + 1) \ 3
\).
\end{example}

Las reglas para evaluar expresiones del \(\lambda\)-cálculo se conocen como
\emph{reglas de reducción} y son:
\begin{itemize}
  \item \(\alpha\)-conversión: Permite renombrar variables ligadas:
    \[
    (\lambda x \mapsto M) \ N \rightsquigarrow M[x \coloneqq N]
    \]
  \item Sustitución: Permite reemplazar todas las ocurrencias \emph{libres} de la
    variable \(x\) en la expresión \(M\) con la expresión \(N\) y se denota:
    \(M[x \coloneqq N]\).
  \item \(\beta\)-reducción: Permite aplicar una función a un argumento,
    se define mediante la sustitución:
    \[
    (\lambda x \mapsto M) \ N \rightsquigarrow M[x \coloneqq N]
    \]
  \item \(\eta\)-conversión: Expresa el concepto de \emph{extensionalidad} de funciones,
    esto es, dos funciones son iguales si y solo si devuelven el mismo resultado
    para todos los argumentos. Por ejemplo cuando \(x\) no es libre en \(f\), 
    podemos aplicar la \(\eta\)-reducción:
    \[
    (\lambda x \mapsto f \ x) \rightsquigarrow f
    \]
    En el sentido contrario podemos aplicar \(\eta\)-expansión a \(f\) para
    obtener nuevamente la expresión anterior:
    \[
    f \rightsquigarrow (\lambda x \mapsto f \ x) \ x
    \]
\end{itemize}

\begin{example}
  Podemos aplicar \(\beta\)-reducción a la expresión del Ejemplo~\ref{ex:lambda_expr}:
  \[
  (\lambda x \mapsto x + 1) \ 3 \rightsquigarrow 3 + 1 
  \]
\end{example}

Las reglas de reducción son \enquote{unidireccionales}, es decir, van de izquierda a
derecha. Viendo el ejemplo anterior uno podría pensar que también \enquote{reducen}
expresiones a versiones \enquote{más simples}: aplicando las reglas una tras otra
llegaríamos a una versión de la expresión en la que las reglas de reducción no
se pueden seguir aplicando, esta es la llamada \emph{forma normal} de una expresión.

\begin{definition}[Forma normal]
  Una expresión está en \emph{forma normal} si no se puede reducir más.
\end{definition}

Lamentablemente, no todas las expresiones tienen una forma normal, como veremos en el
siguiente ejemplo.

\begin{definition}[Combinador \(\Omega\)]\label{def:omega_combinator}
  El combinador \(\Omega\) es la expresión \(\lambda\):
  \[
  \Omega \coloneqq (\lambda x \mapsto x \ x) \ (\lambda x \mapsto x \ x) 
  \]
\end{definition}

Sin adentrarnos en los detalles de la definición de la sintaxis y reducción del
\(\lambda\)-cálculo no podemos dar una demostración \enquote{rigurosa} del siguiente
teorema, pero podemos dar una informal.

\begin{theorem}
  El combinador \(\Omega\) no tiene forma normal.
\end{theorem}

\begin{proof}
  Aplicando \(\beta\)-reducción a la expresión~\ref{def:omega_combinator} obtenemos la
  misma expresión:
  \[
  (\lambda x \mapsto x \ x) \ (\lambda x \mapsto x \ x) \rightsquigarrow
  (\lambda x \mapsto x \ x) \ (\lambda x \mapsto x \ x)
  \]
\end{proof}

Como dijimos antes, los únicos \enquote{tipos de datos} presentes en el
\(\lambda\)-cálculo son las funciones, por lo que si queremos utilizar números,
booleans, o cualquier otro tipo de dato debemos \enquote{codificarlos} como
tal. Veamos una codificación usual de los booleanos, conocida como \emph{codificación de Church}.
En esta codificación los booleanos son funciones que toman dos argumentos y
devuelven uno de ellos. Por ejemplo, podemos definir los literales booleanos como:

\begin{definition}[Booleanos de Church]
  Definimos los literales booleanos como:
  \[
    \texttt{true} \coloneqq \lambda x \mapsto \lambda y \mapsto x \qquad
    \texttt{false} \coloneqq \lambda x \mapsto \lambda y \mapsto y 
  \]
\end{definition}

Ahora podemos definir los operadores lógicos \(\texttt{and}\), \(\texttt{or}\), y
\(\texttt{not}\) como funciones que toman dos booleanos y devuelven un booleano.

\begin{definition}[Operadores lógicos]
  Definimos los operadores lógicos como:
  \[
    \texttt{and} \coloneqq \lambda x \mapsto \lambda y \mapsto x \ y \ x \quad
    \texttt{or} \coloneqq \lambda x \mapsto \lambda y \mapsto x \ x \ y \quad
    \texttt{not} \coloneqq \lambda x \mapsto x \ \texttt{false} \ \texttt{true} 
  \]
\end{definition}

Podemos incluso definir la estructura de control \emph{if-then-else}.

\begin{definition}[If-then-else]
  Definimos la estructura de control \emph{if-then-else} como:
  \[
  \texttt{if-then-else} \coloneqq \lambda b \mapsto \lambda x \mapsto \lambda y \mapsto b \ x \ y 
  \]
\end{definition}

Los naturales también tienen una codificación de Church, pero la omitimos ya que
más adelante veremos una forma más general de codificar tipos de datos.

Como hemos visto, el \(\lambda\)-cálculo es un sistema formal simple pero poderoso
que se puede usar para razonar sobre la computación, de hecho es \emph{Turing
completo}, lo que significa que puede expresar cualquier función computable. Sin embargo, como hemos
visto con el \(\Omega\)-combinador, no es muy \enquote{seguro} de usar.

\subsection{Lambda cálculo tipado}

Para evitar el problema introducido por las computaciones no terminates como el
combinador \(\Omega\) podemos introducir \emph{tipos} al cálculo lambda\footnote{
  En algunos sistemas de tipos el combinador \(\Omega\) puede estar bien tipado, este no es
  el caso en los sistemas de tipos normalmente usados en asistentes de pruebas, como el CIC.}.

Como indicamos en la introducción, los tipos son una forma de clasificar las
expresiones del \(\lambda\)-cálculo, y nos permiten asegurar que las
expresiones están bien tipadas, es decir, que las funciones se aplican a los
argumentos correctos.

Añadir tipos también asegura que expresiones sin sentido, como aplicar un booleano a
un entero no son posibles, y que las expresiones no pueden quedar \enquote{atascadas} en una
computación no terminante.

\subsubsection{Lambda cálculo simplemente tipado}

El \emph{cálculo lambda tipado} más simple es el \emph{cálculo lambda simplemente tipado} o
\emph{STLC}. En el STLC cada variable y cada función tiene un tipo, la
abstracción ahora toma la forma \(\lambda x: \alpha \mapsto M\), donde 
\(\alpha\) es el tipo de la variable \(x\). Los dos tipos obligatorios en el STLC
son el tipo función \(\alpha \to \beta\), que representa una función que
toma un argumento de tipo \(\alpha\) y devuelve un valor de tipo \(\beta\), y
el tipo base \(\texttt{Bool}\), que representa los booleanos.

La sintaxis del STLC es la siguiente:
\begin{align*}
  M \Coloneqq x & \mid \lambda x: \alpha \mapsto M \mid M \ M \\
    & \mid \texttt{true} \mid \texttt{false} \\
    & \LeanIfElse{M}{M}{M}
\end{align*}

Además de las reglas de reducción que vimos antes, debemos añadir las
\emph{reglas de tipado}, que son reglas que definen las expresiones bien
tipadas del cálculo. Estas reglas requieren un \emph{contexto de tipado} que
normalmente se denota como \(\Gamma\), este es un mapa que asocia variables con
tipos. Estas reglas, por ejemplo, aseguran que la estructura \emph{if-then-else} se aplica
a booleanos, y que las dos ramas tienen el mismo tipo.

\subsubsection{Generalización}

El STLC es la base de la mayoría de los \(\lambda\)-cálculos tipados, y se puede
extender permitiendo \emph{términos que dependen de tipos}, esto se conoce como
\emph{polimorfismo} y el cálculo se conoce como \emph{System F}. También podemos
extender el STLC con \emph{tipos que dependen de otros tipos}, que llamaremos
\emph{constructores de tipos}, y el cálculo se conoce como \emph{System F\(_\omega\)}.
Si añadimos \emph{tipos dependientes} obtenemos \(\lambda P\), etc\ldots

Podemos combinar estas extensiones de forma arbitraria para obtener sistemas
con capacidad expresiva distinta, una taxonomía de estos sistemas es el 
\(\lambda\)-cubo de Barendregt. En la cúspide de la potencia expresiva tenemos
el CoC, que es la base de Coq y Lean. En lo que queda del texto exploraremos
distintas características del CoC. Para un tratamiento completo del
\(\lambda\)-cubo referimos al lector a~\cite{lambda_with_types}.

\subsection{Teoría de tipos y teoría de conjuntos}

En esta sección seguiremos la presentación de la teoría de tipos
desarrollada en~\cite{hottbook}, aunque los conceptos presentados son
comunes a todas las teorías de tipos de Martin-Löf. En rigor, no existe 
\emph{una teoría de tipos}, si no varios sistemas formales basados en las mismas ideas.
Las teorías de tipos de Martin-Löf, 
el Cálculo de Construcciones, y el Cálculo de
Construcciones Inductivas son distintas, pero para los propósitos de este texto
nos referiremos a todas ellas como \emph{teoría de tipos}.

Podemos pensar en el \(\lambda\)-cálculo presentado antes como una \enquote{sintaxis} para
nuestra teoría de tipos, los tipos le darán significado a las expresiones del
\(\lambda\)-cálculo. De hecho, una buena parte de las secciones siguientes estarán
dedicadas a mostrar cómo \emph{construir} tipos que pueden codificar los objetos
de nuestras pruebas.

Una diferencia fundamental entre tipos y conjuntos es que los tipos no
\emph{contienen} valores, sino que \emph{clasifican} valores. En teoría de conjuntos un
valor es un \emph{miembro} de un conjunto, mientras que en teoría de tipos un
valor \emph{tiene} un tipo. La notación análoga a la de pertenencia en teoría de
conjuntos (\(x \in S\)) es el \emph{juicio de tipado} (\emph{typing judgement}) en teoría de tipos, 
que se escribe como \(x: \alpha\), y se lee como \enquote{x tiene tipo
\(\alpha\)}.
Mientras que en teoría de tipos \(x \in S\) y \(x \notin S\) son expresiones
(en el sentido lógico) válidas, el juicio de tipado es \enquote{siempre verdadero}, no es
una proposición que se pueda probar o refutar. Debido a esto no podemos
introducir una variable sin especificar su tipo, cuando escribimos \(x\) estamos
implícitamente escribiendo \(x: \alpha\), donde \(\alpha\) es el tipo de \(x\).

Otra importante diferencia es que la teoría de tipos es \emph{su propio sistema
deductivo}, mientras que la teoría de conjuntos está formulada dentro de la
lógica de primer orden. Esto significa que la teoría de conjuntos tiene dos
nociones primitivas, conjuntos y proposiciones, mientras que la teoría de tipos
tiene tipos y funciones. Con solo estas dos podemos recrear toda la lógica de primer
orden, y de hecho, podemos usar tipos para codificar proposiciones, y funciones
para codificar demostraciones. El juicio de tipo \(x: \alpha\) también puede ser leído como
\enquote{x es una demostración de \(\alpha\)}, y el tipo de función \(\alpha \to \beta\) puede ser leído como
\enquote{si tenemos una demostración de \(\alpha\) podemos construir una demostración de
\(\beta\)}. Esto significa que la abstracción de funciones corresponde a una
\emph{implicación lógica} en la lógica clásica, y la aplicación de funciones
corresponde a la \emph{eliminación de implicaciones}. Veremos más sobre esto en la
Sección~\ref{sec:curry_howard}.

\subsection{Igualdad}

La igualdad en teoría de tipos es ligeramente más compleja que en las matemáticas
clásicas. La igualdad que tenemos \enquote{gratis} es la \emph{igualdad
definicional}, es decir, dos términos son iguales si pueden ser transformados
en otros mediante \emph{reglas de reescritura} (análogas a las reglas de reducción
del \(\lambda\)-cálculo). Las reglas de reducción de las teoría de tipos
típicamente incluyen:
\begin{itemize}
  \item \(\beta\)-reducción. Reemplazar la aplicación de una función por el
    cuerpo de la función, donde el argumento se reemplaza por la variable ligada.
  \item \(\delta\)-reducción. Reemplazar una constante por su definición.
  \item \(\iota\)-reducción. Simplificar expresiones con constructores y encaje
    de patrones para tipos de datos inductivos.
\end{itemize}

Estas reglas satisfacen normalmente un número de propiedades deseables, como
\emph{confluencia}, \emph{terminación}, y \emph{consistencia}, que aseguran que el
proceso de reducción está bien definido, y que la relación de igualdad es una
relación de equivalencia. En algunos casos se añaden otras propiedades como
\emph{congruencia}, \emph{normalización fuerte} y \emph{canonicidad}.

Otro tipo de igualdad es la \emph{igualdad proposicional}, que establece que dos
términos son iguales si se puede demostrar que son iguales en la teoría de
tipos. Esta es más parecida a la noción clásica de igualdad, y se usa para
probar propiedades de los términos. Un caso especial de igualdad proposicional es
la \emph{extensionalidad funcional}, que establece que dos funciones son
iguales si son iguales para todos los argumentos, un resultado clásico de la
teoría de conjuntos que sigue del axioma de extensionalidad, que establece que
dos conjuntos son iguales si contienen los mismos elementos. (La relación puede
depender del axioma K).

Otro axioma importante es la \emph{extensionalidad proposicional}, que
establece que dos proposiciones son iguales si son equivalentes, es decir, si se
implican mutuamente. Esta es una propiedad clave de la teoría de tipos, ya que
permite reescribir proposiciones en las pruebas y simplificar el razonamiento.

La igualdad es también una \emph{relación binaria}, es decir, una relación que relaciona
dos valores. En teoría de tipos las relaciones binarias son funciones
que toman dos valores y devuelven una proposición, esta proposición es verdadera
si la relación se cumple, y falsa en caso contrario.

\begin{definition}[Relación binaria]
  Una relación binaria es una función que toma dos elementos y devuelve una
  proposición.
  \[
  \texttt{Rel} \ (\alpha : \texttt{Type}) \ (\beta : \texttt{Type}) 
    \coloneqq \alpha \to \beta \to \Prop 
  \]
\end{definition}

Podemos entonces definir la relación de igualdad como:

\begin{definition}[Relación de igualdad]
  La relación de igualdad es una relación binaria que toma dos elementos de un
  mismo tipo y devuelve una proposición.
  \[
  \texttt{Eq} \ (\alpha: \texttt{Type}) \coloneqq \texttt{Rel} \ \alpha \ \alpha  
  \]
\end{definition}

La igualdad es, sin embargo, mucho más que una relación de equivalencia. Tiene la
importante propiedad de que cada afirmación respeta la equivalencia, en el sentido
de que podemos sustituir expresiones iguales sin cambiar el valor de verdad.

\subsubsection{Equivalencia}

La noción de equivalencia es más general que la igualdad, y se refiere a una relación
entre dos tipos que es reflexiva, simétrica y transitiva. Esta relación junto
con el tipo en el que se define conforma un \emph{setoid}, que es análogo 
a los \emph{conjuntos cociente} en teoría de conjuntos. Es decir nos define
\emph{isomorfismos} entre tipos.

\subsection{Universos}

Al principio de la sección hemos introducido el juicio \(x: \alpha\), en el que
\(x\) es un \enquote{valor} y \(\alpha\) es un \enquote{tipo}, pero ahora nos preguntamos
¿cuál es el tipo de \(\alpha\)? La respuesta es que \(\alpha\) es de tipo
\(\texttt{Type}\), pero ¿cuál es el tipo de \(\texttt{Type}\)? Por supuesto también es
un tipo, así que \(\texttt{Type}: \texttt{Type}\). 
Esto lleva a un problema conocido como \emph{impredicatividad}.

La impredicatividad surge cuando una definición de un objeto o tipo depende de un
cuantificador que abarca una totalidad que incluye el propio objeto que se está
definiendo. Este aspecto autorreferencial puede llevar a paradojas lógicas,
como la paradoja de Russell en teoría de conjuntos, donde un conjunto se define
como aquel que contiene todos los conjuntos que no se contienen a sí mismos.

El análogo a la paradoja de Russell en teoría de tipos es la \emph{paradoja de
Girard}, en un sistema que permite la impredicatividad, podemos definir un tipo \(U\)
que \enquote{contiene} todos los tipos que no se \enquote{contienen} a sí mismos, y luego preguntar si
\(U: U\), si lo es, entonces no debería ser de tipo \(U\), y si no lo es, entonces
debería serlo. Un sistema que muestra esta paradoja no es muy útil para la
demostración de teoremas, ya que es inconsistente, por lo que cualquier
proposición puede ser probada.\footnote{
  Históricamente este problema se ha resuelto con el \emph{Sistema U}.} Para evitar
esto, introducimos una jerarquía de \emph{universos de tipos}.

Cada universo es una \emph{colección de tipos}\footnote{Si estos son tipos o no depende de la
  teoría de tipos en particular, en Lean y Coq no lo son}, cada universo contiene los
tipos del universo anterior, y está contenido en el siguiente universo. Por
ejemplo: \(\texttt{Type}: \texttt{Type} \ 1\),
\(\texttt{Type} \ 1: \texttt{Type} \ 2\), y así sucesivamente. El tipo función
habita el universo más pequeño que contiene los tipos de su dominio y codominio.
Por ejemplo, \(\mathbb{N} \to \mathbb{N}: \texttt{Type} \ 1\), y
\(\texttt{Type} \ 1 \to \texttt{Type} \ 2: \texttt{Type} \ 3\).

La base de la jerarquía, el universo \(\texttt{Type} \ 0\), contiene los tipos del
lenguaje de programación en sí, como \(\mathbb{N}\), \(\texttt{Bool}\), y
\(\texttt{List}\). El segundo universo, \(\texttt{Type} \ 1\), contiene los tipos
del primer universo, y así sucesivamente. Esta jerarquía de universos evita la
impredicatividad, ya que cada universo es un tipo distinto y no puede contenerse
a sí mismo.

\subsection{Constructores}

Hemos visto varias propiedades de los tipos, pero no hemos visto cómo crear
valores de estos tipos. Por ejemplo, para el tipo \(\texttt{Bool}\), ¿cómo
podemos crear el valor \(\texttt{true}\)? Para esto necesitamos \emph{constructores de
tipos}, que son funciones que crean elementos de un tipo y definen las
posibles formas que un elemento de un tipo puede tomar. Para \(\texttt{Bool}\) tenemos dos
constructores, \(\texttt{true}\) y \(\texttt{false}\), que crean los elementos del
tipo. Esto asegura que los elementos del tipo son \emph{solo} \(\texttt{true}\) o
\(\texttt{false}\), y no otro valor.

\subsection{Encaje de patrones}

Para definir funciones que operan sobre tipos necesitamos una forma de distinguir
entre las distintas formas que un elemento de un tipo puede tomar, esto se
hace usando \emph{encaje de patrones}. El encaje de patrones es una forma de
definir funciones definiendo un caso distinto para cada constructor del tipo. Por
ejemplo, podemos definir una función \(\texttt{not}\) que toma un booleano y
devuelve su negación como:
\[
  \texttt{not} \ \texttt{true} = \texttt{false} \qquad
  \texttt{not} \ \texttt{false} = \texttt{true}
\]

\subsection{Tipos suma y producto}

Hemos visto que podemos construir tipos a partir de otros mediante funciones,
otro mecanismo que tenemos para construir tipos son los \emph{tipos suma} y
\emph{tipos producto}. Estos tipos son construcciones primitivas de la
teoría de tipos, y son usados para construir tipos complejos a partir de tipos
más simples. Los tipos suma son equivalentes a la \emph{unión disjunta} en teoría de
conjuntos, y son usados para representar tipos que pueden tomar uno de varios
valores. Denotamos un tipo suma como \(\alpha + \beta\), donde \(\alpha\) y \(\beta\) son
tipos, los dos constructores de este tipo son \(\texttt{inl}\) y
\(\texttt{inr}\) (abreviaciones de \emph{inyección izquierda} e
  \emph{inyección derecha}), que crean un valor del tipo suma. 
Mediante encaje de patrones podemos definir funciones que operan sobre tipos suma,
distinguiendo entre los dos constructores. 

Cuando hemos definido antes el tipo \(\texttt{Bool}\) como
el tipo que puede tomar un valor de \(\texttt{true}\) o \(\texttt{false}\), hemos
utilizado implícitamente un tipo suma. En efecto, sea \(\texttt{Unit}\) el tipo
\emph{unitario} (que tiene un único constructor \texttt{()}), podemos definir 
\(\texttt{Bool}\) como el \emph{tipo binario} que es la suma de \(\texttt{Unit}\) y
\(\texttt{Unit}\):
\[
  \texttt{Bool} \coloneqq \texttt{Unit} + \texttt{Unit}
\]
Donde la inyección izquierda \(\texttt{inl}\) representa \(\texttt{true}\) y la
inyección derecha \(\texttt{inr}\) representa \(\texttt{false}\).
Mediante encaje de patrones podríamos definir la función \texttt{not} como:
\begin{align*}
  \texttt{not} \ (\texttt{inl} \ \placeholder) &= \texttt{inr} \ \texttt{()} \\
  \texttt{not} \ (\texttt{inr} \ \placeholder) &= \texttt{inl} \ \texttt{()}
\end{align*}

\begin{notation}
  Utilizamos el símbolo \(\placeholder\) (\emph{placeholder} o \emph{marcador de
  posición} en inglés) para denotar un argumento que no
  utilizamos en una expresión, o que puede ser inferido por las reglas de tipado.
\end{notation}

Y la función \texttt{and} como:
\begin{align*}
  \texttt{and} \ (\texttt{inr} \ \placeholder) \ \placeholder &= \texttt{inr} \ \texttt{()} \\
  \texttt{and} \ (\texttt{inl} \ \placeholder) \ y &= y
\end{align*}
Esto, aunque interesante en teoría, no es muy practico, enseguida veremos una
generalización de los tipos suma, los \emph{tipos inductivos}.

Los tipos producto son equivalentes al \emph{producto cartesiano} en teoría de
conjuntos, y son usados para representar tipos que pueden tomar varios valores
al mismo tiempo, visto de otra forma \enquote{parejas} de valores. Denotamos un tipo
producto como \(\alpha \times \beta\), donde \(\alpha\) y \(\beta\) son tipos, su
único constructor es \(\texttt{intro}: \alpha \to \beta \to \alpha \times \beta\), 
que crea un valor del tipo producto.
Definimos también las \emph{proyecciones} \(\texttt{fst}\) y
\(\texttt{snd}\), que devuelven el primer y segundo elemento de la pareja
respectivamente.

Podemos definir por ejemplo el tipo \(\texttt{And}\) (que no se debe confundir
con la función \texttt{and} de \texttt{Bool}) como el tipo producto de dos 
proposiciones\footnote{Veremos el sentido de esto en la sección~\ref{sec:curry_howard}}:
\[
  \texttt{And} \coloneqq \Prop \times \Prop
\]
De esta forma la demostración izquierda de la conjunción 
\(x : \texttt{And} \coloneqq (p, q)\) es \(\texttt{fst} \ x = p\)
y la derecha \(\texttt{snd} \ x = q\). La generalización a n-tuplas de
valores es inmediata.

\subsection{Tipos inductivos}

Definir tipos finitos como \(\texttt{Bool}\) es interesante, pero no nos
llevará muy lejos, necesitamos una forma de definir tipos infinitos, como los
números naturales. Esto se hace usando \emph{tipos inductivos}, que son tipos
definidos por un conjunto de constructores que definen las posibles formas que
un elemento del tipo puede tomar. Por ejemplo, podemos definir los números
naturales como un tipo inductivo con dos constructores, \(0\) y
\(\texttt{succ}\), que crean los elementos del tipo, de una forma análoga a los
\emph{axiomas de Peano}. En este caso los números naturales tienen los
constructores:
\[
  0 : \mathbb{N} \qquad \texttt{succ}: \mathbb{N} \to \mathbb{N}
\]

De forma similar al principio clásico de inducción, hemos definido un caso base
\(0\) y un caso inductivo \(\texttt{succ}\), que crea el sucesor de un número
natural. De esta forma podemos crear los números naturales como un tipo
inductivo, y definir funciones que operan sobre ellos, como la suma,
multiplicación y exponenciación. Esta es la esencia de los tipos inductivos, la
capacidad de definir tipos que son creados por un conjunto de constructores, y
definir funciones que operan sobre ellos. Y como es el caso en la matemática
clásica, podemos usar la inducción para probar propiedades de estos tipos, en
este caso \emph{inducción estructural}.

Siguiendo el ejemplo de los números naturales, podemos construir una función
\(f: \mathbb{N} \to C\) por recursión sobre los naturales, es suficiente con
proveer un caso base \(c_0 : C\) y una función de \enquote{siguiente paso}
\(c_s: \mathbb{N} \to C \to C\), luego podemos definir \(f\) mediante encaje de
patrones como:
\begin{align*}
  f \ 0 & = c_0 \\
  f \ (\texttt{succ} \ n) & = c_s \ n \ (f \ n)
\end{align*}
Para usar esta técnica debemos asegurar que el tipo inductivo es \emph{bien fundado},
es decir, que los constructores del tipo son aplicados un número finito de veces
para crear un elemento del tipo.

\begin{example}[Suma de naturales]
  Definimos la suma de dos números naturales como:
  \begin{align*}
    \texttt{add} \ n \ 0 &= n \\
    \texttt{add} \ n \ (\texttt{succ} \ m) &= \texttt{succ} \ (\texttt{add} \ n \ m)
  \end{align*}
\end{example}

\subsection{Tipo función dependiente}

Como dijimos antes, las funciones son ciudadanos de primera clase en la teoría de tipos,
y pueden ser pasadas como argumentos, devueltas como resultados, y almacenadas en
variables. Nos proponemos generalizar el tipo función de forma que el tipo de
retorno dependa del valor de los argumentos, llamaremos a esta generalización \emph{\(\Pi\)-tipos}
(o \emph{tipos dependientes}).

La notación para un \(\Pi\)-tipo es \(\Pi_{x: \alpha} B \ x\), donde \(x\) es el argumento
del cual depende el tipo, \(\alpha\) es el tipo del argumento, y \(B\) es una familia
de tipos indexados por \(x\) (\(B: \alpha \to \texttt{Type}\)). Por ejemplo, podemos
definir una función que toma un número natural \(n\) y devuelve un tipo
\(\texttt{Vec} \ \alpha \ n\), que es un tipo de vector de longitud \(n\) de
elementos de tipo \(\alpha\).
\begin{example}[Vector]
  % We can define a vector of a given length as an inductive type that depends on
  % the length of the vector. We define the type \(\texttt{Vec}\) as:
  Definimos un vector de longitud \(n\) como un tipo inductivo que depende de la
  longitud del vector
  \[
  \BigPi_{n: \mathbb{N}} \texttt{Vec} \ \alpha \ n
  \]
  Con constructores:
  \begin{align*}
    \texttt{nil} & : \texttt{Vec} \ \alpha \ 0 \\
    \texttt{cons} & : \BigPi_{n: \mathbb{N}}
    (\alpha \to \texttt{Vec} \ \alpha \ n
      \to \texttt{Vec} \ \alpha \ (\texttt{succ} \ n)) 
  \end{align*}
\end{example}

De esta forma los \(\Pi\)-tipos nos proporcionan una forma de tipar estructuras de
datos complejas, como listas, árboles y grafos cuyo tipo depende del valor
de sus elementos. 

\subsection{Tipo par dependiente}

Ahora nos proponemos generalizar los tipos producto de forma que el tipo de
la segunda componente dependa del valor de la primera, llamaremos a esta
generalización \emph{\(\Sigma\)-tipos} (o \emph{tipos pares dependientes}). La
notación para un \(\Sigma\)-tipo es \(\Sigma_{x: \alpha} B \ x\), donde \(x\) es 
la primera componente del par, \(\alpha\) su tipo, y \(B\) es una familia de
tipos indexados por \(x\) (\(B: \alpha \to \texttt{Type}\)).

En caso de que \(B\) sea una función constante el \(\Sigma\)-tipo es
igual al tipo producto, es decir, si \(B \ x = \beta\) para todo \(x\), entonces:
\[
  \BigSigma_{x: \alpha} B = \alpha \times \beta.
\]

Su constructor es \(\texttt{intro}: \alpha \to B \ x \to \Sigma_{x: \alpha} B \ x\), que toma un
valor de tipo \(\alpha\) y un valor de tipo \(B \ x\) y devuelve un valor del
tipo \(\Sigma_{x: \alpha} B \ x\).

Aunque la utilidad de estos tipos no es tan evidente como la de los \(\Pi\)-tipos, 
veremos en la siguiente sección que son fundamentales.

\subsection{La correspondencia de Curry-Howard}\label{sec:curry_howard}

El isomorfismo de Curry-Howard establece una correspondencia entre \emph{proposiciones}
y \emph{tipos} y entre \emph{términos} y \emph{demostraciones}. Este es conocido como el
\emph{principio de proposiciones como tipos}, y es la base de todos los
asistentes de prueba modernos, como Coq, Lean y Agda. En nuestro sistema de tipos
una proposición es de tipo \(\Prop\), las \enquote{demostraciones} son términos de
ese tipo, y la proposición es verdadera si el tipo está habitado, es decir, si
hay un término de ese tipo, este es nuestro \emph{testigo}. 
Consideremos por ejemplo el teorema \enquote{\(\sqrt{2}\) es irracional}, 
esta proposición vive en el tipo \(\Prop\), y su demostración es un
término de ese tipo, que sabemos que está habitado.

Las demostraciones pueden ser manipuladas como en la lógica clásica, por ejemplo
podemos obtener una demostración de que dos proposiciones son ciertas usando el
operador \(\land\), que establece que si \(P\) y \(Q\) son verdaderas, entonces
\(P \land Q\) es verdadero. Como hemos visto antes esto corresponde al \emph{tipo
producto} \(\texttt{And} \coloneqq P \times Q\), que establece que si tenemos una 
demostración de \(P\) y una demostración de \(Q\), entonces podemos construir una
demostración de \(P \times Q\). El juicio \(\bot\) que representa la
\emph{inconsistencia} se representa como un tipo sin constructores \(\texttt{Empty}\), 
y el juicio \(\top\) que representa una \emph{tautología} se representa como el tipo unitario
\(\texttt{Unit}\), que tiene solo un constructor sin argumentos.

El principio \emph{exfalso quodlibet} se representa como una función de tipo
\(\texttt{exfalso}: \texttt{Empty} \to \alpha\), que establece que si tenemos una
demostración de \(\bot\), entonces podemos construir una demostración de cualquier
proposición \(\alpha\). Claramente, si intentamos definir esta función mediante
encaje de patrones tenemos que no podemos discriminar el valor de retorno
dependiendo del valor de entrada, ya que no hay ningún constructor de
\(\texttt{Empty}\).

Consideramos ahora la implicación (\(\implies\)) en lógica proposicional, la
proposición \(P \implies Q\) establece que si \(P\) es verdadero, entonces \(Q\)
debe ser verdadero. En teoría de tipos esto corresponde al tipo función
\(P \to Q\), que establece que si tenemos una demostración de \(P\), entonces
podemos construir una demostración de \(Q\). La demostración de la proposición
corresponde a un término del tipo, y la demostración de la implicación
corresponde a una función que toma una demostración de \(P\) y devuelve una
demostración de \(Q\).

La correspondencia no se limita a la lógica proposicional, sino que se
extiende a la lógica de predicados y a lógicas de orden superior. Por ejemplo, el
cuantificador universal \(\forall\) corresponde al tipo función dependiente
\(\Pi\), en efecto, la proposición \(\forall x: \alpha, P \ x\) corresponde al tipo
\(\Pi_{x: \alpha} P \ x\), y el cuantificador existencial \(\exists x : \alpha, P \ x\)
corresponde al tipo par dependiente \(\Sigma_{x: \alpha} P \ x\). Donde las 
\emph{familias indexadas} \(P: \alpha \to \Prop\) son las proposiciones 
que dependen de un valor. De esta forma podemos razonar sobre programas usando las
herramientas de la lógica, y viceversa.

Clarificamos esto con un ejemplo, la desigualdad entre naturales \(n \leq m\) se
puede expresar como la proposición \(n \leq m \coloneqq \exists k: \mathbb{N}, n + k = m\).
Para demostrar entonces que \(3 \leq 5\) debemos construir un término del tipo
\(\Sigma_{k: \mathbb{N}}, \texttt{add} \ 3 \ k = 5\). Para ello usamos el constructor
de \(\Sigma\)-tipos:
\[
  \texttt{intro} \ 2 \ (\texttt{add} \ 3 \ 2 = 5)
\]
Ahora podemos \enquote{evaluar} la función \(\texttt{add}\) obteniendo:
\[
  \texttt{intro} \ 2 \ (5 = 5)
\]
y aplicando el constructor de igualdad \(\texttt{refl}\):
\[
  \texttt{intro} \ 2 \ (\texttt{refl} \ 5)
\]
conseguimos un término del tipo \(\Sigma_{k: \mathbb{N}} \texttt{add} \ 3 \ k = 5\), que
es una demostración de \(3 \leq 5\). Este ejemplo muestra como realizaremos 
todas las demostraciones en Lean a partir de ahora.

El \emph{verificador de tipos} del sistema nos permite entonces \emph{verificar} las
demostraciones, comprobando que los términos están bien tipados, y que los términos
corresponden a demostraciones válidas. Este mecanismo asegura que las
demostraciones son correctas, y que los programas están bien tipados.

La \enquote{demostración} de que la teoría de tipos y la teoría de conjuntos son
\enquote{igual de potentes} consiste en recrear la teoría de conjuntos en la teoría de
tipos, y viceversa. Esto va más allá de los objetivos de este trabajo, pero
en la siguiente sección veremos como podemos recrear conjuntos dentro
de la teoría de tipos, ya que más adelante los necesitaremos.

\subsection{Conjuntos}\label{sec:sets}

Podemos codificar un conjunto en teoría de tipos como una función que toma un
valor de un tipo y devuelve una proposición, esta proposición nos dice si el
valor está en el conjunto o no.

\begin{definition}[Conjunto]
  Un conjunto es una función que toma un elemento de un tipo y devuelve una
  proposición.
  \[
  \texttt{Set} \ (\alpha: \texttt{Type}) \coloneqq \alpha \to \Prop 
  \]
\end{definition}

Podemos usar también la notación tradicional de \emph{generador de conjuntos}.
\begin{notation}[Generador de conjuntos]
  Definimos un conjunto usando la notación de generador de conjuntos como:
  \[
  \{x \mid P\} \coloneqq \lambda x \mapsto P 
  \]
\end{notation}

La definición de la relación de pertenencia es inmediata, ya que la
propia definición del conjunto captura el concepto de pertenencia.
\begin{definition}[Pertenencia]
  La relación de pertenencia es una función que toma un valor de un tipo y un
  conjunto y devuelve una proposición, sea \(x: \alpha\) y \(S: \texttt{Set} \ \alpha\):
  \[
  x \in S \coloneqq S \ x 
  \]
\end{definition}

Los subconjuntos se definen como en teoría de conjuntos.
\begin{definition}[Subconjunto]
  Un subconjunto es una función que toma dos conjuntos y devuelve una
  proposición, sea \(A \ B: \texttt{Set} \ \alpha\):
  \[
  A \subseteq B \coloneqq \forall x, x \in A \implies x \in B 
  \]
\end{definition}

\begin{definition}[Subconjunto estricto]
  Un subconjunto estricto es un subconjunto diferente del conjunto en sí, sea
  \(A \ B: \texttt{Set} \ \alpha\), entonces:
  \[
  A \subset B \coloneqq A \subseteq B \land A \neq B
  \]
\end{definition}

Con estas definiciones las operaciones conjuntistas se definen de la forma
habitual, y podemos razonar sobre ellas usando las herramientas de la teoría de
tipos.
\begin{definition}[Operaciones conjuntistas]
  Definimos las operaciones conjuntistas como:
  \begin{align*}
    A \cup B & \coloneqq \{x \mid x \in A \lor x \in B\} \\
    A \cap B & \coloneqq \{x \mid x \in A \land x \in B\} \\
    A \setminus B & \coloneqq \{x \mid x \in A \land x \notin B\} \\
    A^\complement & \coloneqq \{x \mid x \notin A\} 
  \end{align*}
\end{definition}

Necesitaremos usar conjuntos cuando hagamos nuestra introducción a la
\emph{teoría de dominios}.

\subsection{Irrelevancia de demostraciones}

Una diferencia fundamental entre las diferentes teorías de tipos es cómo
tratan la igualdad entre demostraciones. El sistema de tipos de Lean, como hemos
dicho antes, tiene un universo especial \(\Prop\), en la base de su jerarquía
de tipos, que contiene las proposiciones. En este universo, todas las demostraciones
de la misma proposición son \emph{iguales}.

Esto quiere decir que si tenemos, por ejemplo, dos demostraciones de que \(2 + 2 = 4\), una
que usa los axiomas de Peano, y otra que usa la definición de suma, estas dos
demostraciones son \emph{iguales}, y pueden ser usadas indistintamente.
Más formalmente:
\[
  \forall P: \Prop, \forall p_1, p_2: P, p_1 = p_2
\]

La ventaja de esto es que nos permite simplificar demostraciones complejas,
ya que no tenemos que preocuparnos por la forma en que se construyen las
demostraciones, sino que podemos tratarlas como \enquote{iguales}. Esto también
ayuda a la hora de \emph{extraer código} de las demostraciones, ya que podemos
\enquote{borrar} las demostraciones en tiempo de compilación y quedarnos solo con 
el resultado. Podemos, por ejemplo, crear una estructura de datos como una lista,
definir sus operaciones y demostrar sus propiedades, y, una vez que estemos
seguros de su corrección, extraer el código a otro lenguaje, como OCaml,
sin incurrir en el coste en tiempo de ejecución de ejecutar las demostraciones.

El axioma de \emph{irrelevancia} de demostraciones implica otro axioma conocido
como \emph{unicidad de las demostraciones de identidad} o \emph{UIP}.
\begin{definition}[UIP]
  Sean \(x \ y : \alpha\) y \(p \ q : x = y\), entonces \(p = q\).
\end{definition}
Esto se debe a que en Lean \(\texttt{Eq} \ \alpha \ x \ y\) tiene tipo \(\Prop\), y
por lo tanto, todas las demostraciones de igualdad son iguales.

Como curiosidad, o potencial desventaja, de este enfoque notamos que UIP no es
compatible con el \emph{axioma de univalencia} de la teoría de tipos homotópica.

En Coq y Agda, por el contrario, no tenemos irrelevancia de demostraciones.
Esta es una de las mayores diferencias entre estos demostradores.

\section{Semántica de paso largo}

% Now that we have introduced type theory and proof assistants, we can move on to
% using them to define the semantics of programming languages. The exposition
% will follow closely the one in~\cite{winskel_formal_semantics}.
% Semantics are the
% study of the \emph{meaning} of programs, and can be defined in different ways, we
% will focus on operational semantics, and denotational semantics.
Una vez presentada la teoría de tipos (y los asistentes de demostración) podemos 
pasar a definir la semántica de lenguajes de programación. Seguiremos un desarrollo
similar al de~\cite{winskel_formal_semantics}. La semántica es el estudio del
\emph{significado} de los programas, y puede ser definida de diferentes formas,
en este trabajo nos centraremos en la semántica operacional y la semántica
denotacional.

Para nuestro estudio primero definiremos un pequeño lenguaje imperativo
\(\texttt{IMP}\) que usaremos para ilustrar los conceptos. \(\texttt{IMP}\) es un
que consiste en un conjunto de comandos (como asignaciones, secuencias,
condicionales y bucles) y expresiones (aritméticas y booleanas). El lenguaje es
sencillo por diseño, para que podamos centrarnos en introducir las técnicas de
la semántica formal, pero estas mismas técnicas pueden ser extendidas a lenguajes
de programación ampliamente usados, como C, Java y Python.

\subsection{Sintaxis}

Como introducción informal a la sintaxis de \(\texttt{IMP}\) podemos definir los
\enquote{conjuntos sintácticos} del lenguaje como:
\begin{itemize}
  \item Enteros \(n \in \mathbb{Z}\).
  \item Booleanos \(b \in \texttt{Bool} = \{\texttt{true}, \texttt{false}\}\).
  \item Variables \(x \in \texttt{Var}\), that store integers.
  \item Expresiones aritméticas \(a \in \texttt{Aexp}\), que pueden ser constantes,
    variables u operaciones aritméticas.
  \item Expresiones booleanas \(b \in \texttt{Bexp}\), que pueden ser constantes,
    variables u operaciones booleanas.
  \item Comandos \(c \in \texttt{Com}\), que pueden ser asignaciones, condicionales,
    bucles o secuencias de comandos.
\end{itemize}

Definimos la sintaxis de \(\texttt{Aexp}\), en forma de Backus-Naur:
\[
a \Coloneqq n \ | \ x \ | \ a_1 \ImpPlus a_2 \ | \ a_1 \ImpMinus a_2 \ | \ a_1 \ImpTimes a_2
\]
la sintaxis de \(\texttt{Bexp}\) como:
\[
b \Coloneqq \texttt{true} \ | \ \texttt{false} \ | \ 
  \ImpNot b \ | \ b_1 \ \ImpAnd \ b_2 \ | \ b_1 \ \ImpOr \ b_2 \ | \ a_1 \ \ImpEq \ a_2 \ | \ a_1 \ \ImpLe \ a_2
\]
y la sintaxis de \(\texttt{Com}\) como:
\[
c \Coloneqq \ImpSkip \ | \ \ImpAssign{x}{a} \ | \ \ImpConcat{c_1}{c_2} \ | \ \ImpIfElse{b}{c_1}{c_2} \ | \ \ImpWhile{b}{c}
\]

Armados con nuestra teoría de tipos podemos definir estos \enquote{conjuntos sintácticos}
como tipos inductivos, el tipo de expresiones aritméticas puede ser definido como un
tipo inductivo \(\texttt{Aexp}\) con constructores para cada \enquote{constructo
sintáctico}:
\begin{align*}
  \texttt{val} &\coloneqq \mathbb{Z} \to \texttt{Aexp} &
  \texttt{var} &\coloneqq \texttt{String} \to \texttt{Aexp} \\
  \texttt{add} &\coloneqq \texttt{Aexp} \to \texttt{Aexp} \to \texttt{Aexp} &
  \texttt{sub} &\coloneqq \texttt{Aexp} \to \texttt{Aexp} \to \texttt{Aexp} &
  \texttt{mul} &\coloneqq \texttt{Aexp} \to \texttt{Aexp} \to \texttt{Aexp}
\end{align*}

Mientras que los primeros dos simplemente \enquote{coaccionan} un valor o una
variable en una expresión, los otros \enquote{construyen} una expresión a partir de
dos expresiones. Los constructores para expresiones booleanas son definidos de
forma similar:
\begin{align*}
  \texttt{true}  &\coloneqq \texttt{Bexp} &
  \texttt{false} &\coloneqq \texttt{Bexp} \\
  \texttt{not}   &\coloneqq \texttt{Bexp} \to \texttt{Bexp} &
  \texttt{and}   &\coloneqq \texttt{Bexp} \to \texttt{Bexp} \to \texttt{Bexp} &
  \texttt{or}    &\coloneqq \texttt{Bexp} \to \texttt{Bexp} \to \texttt{Bexp} \\
  \texttt{eq}    &\coloneqq \texttt{Aexp} \to \texttt{Aexp} \to \texttt{Bexp} &
  \texttt{le}    &\coloneqq \texttt{Aexp} \to \texttt{Aexp} \to \texttt{Bexp}
\end{align*}

Y los comandos son definidos como:
\begin{align*}
  \ImpSkip &\coloneqq \texttt{Com} &
  \texttt{assign} &\coloneqq \texttt{String} \to \texttt{Aexp} \to \texttt{Com} &
  \texttt{seq} &\coloneqq \texttt{Com} \to \texttt{Com} \to \texttt{Com} \\
  \texttt{ifElse} &\coloneqq \texttt{Bexp} \to \texttt{Com} \to \texttt{Com} \to \texttt{Com} &
  \texttt{whileLoop} &\coloneqq \texttt{Bexp} \to \texttt{Com} \to \texttt{Com}
\end{align*}

Una expresión simple en nuestro lenguaje podría ser:
\[
  \texttt{add} \ (\texttt{var} \ \textit{``x''}) \ (\texttt{val} \ 3)
\]
que corresponde al \enquote{árbol sintáctico abstracto} de la expresión, para
no tener que lidiar constantemente con ASTs utilizaremos algunos \enquote{abusos de
notación}, que corresponden al parseo, y diremos que \(x + 3\), por ejemplo, 
es una expresión aritmética de \texttt{IMP}.
Cuando definamos nuestro lenguaje en Lean debemos tener más cuidado al mezclar
nuestro lenguaje y el meta-lenguaje, usando algunas facilidades que nos
proporciona podemos hacer un pequeño parser para \texttt{IMP}, y usando
clases y coacciones podemos incluso usar nuestros abusos de notación en el
lenguaje huésped, todos los detalles están en \url{
github.com/remind-me-later/SemanticsLean/blob/main/Semantics/Imp/Lang.lean}.

Con nuestros abusos un poco más claros, podemos escribir en \(\texttt{IMP}\)
un programa que calcule el factorial de un número \(n\).
\begin{example}[Factorial en \texttt{IMP}]
  El factorial de un número \(n\) puede ser calculado usando un simple bucle:
  \begin{minted}{rust}
  n = 5; i = 2; res = 1;

  while i <= n {
      res = res * i;
      i = i + 1
  }
  \end{minted}
\end{example}

\subsection{Estado}

Como hemos visto en la sintaxis, nuestro lenguaje imperativo soporta una
característica que caracteriza a todos los lenguajes imperativos, valga la redundancia, 
tenemos variables, que \emph{conservan estado} y asignaciones que \emph{modifican el estado}.
La definición de estado es más sutil de lo que uno podría esperar a primera
vista. Un estado es una función que mapea variables a valores, lo podríamos
definir como:
\[
\funfont{state} : \tyfont{String} \to \mathbb{Z}
\]

Pero vamos a definirlo de forma más general, como un \emph{mapa}.
\begin{definition}[Mapa]
  Un mapa es una función que mapea claves (cadenas) a valores:
  \[
  \tyfont{Map} \ (\alpha : \tyfont{Type}) \coloneqq \tyfont{String} \to \alpha
  \]
\end{definition}

Aunque la definición del tipo mapa no es particularmente útil si no tenemos una forma de
relacionar claves y valores, esto se hace usando la función \funfont{update}:
\begin{definition}[Actualizar]
  La función \funfont{update} toma un mapa, una clave y un valor, y devuelve un
  un nuevo mapa que mapea la clave al valor, y todas las otras claves al mismo
  valor que el mapa original.
  \begin{align*}
  \funfont{update} \ & (m: \tyfont{Map} \ \alpha) \ (k: \tyfont{String}) \ (v: \alpha): \tyfont{Map} \ \alpha \coloneqq \\
                      &\lambda k' \mapsto \funfont{cond} \ (k = k') \ v \ (m \ k')
  \end{align*}
\end{definition}

La función \funfont{cond} no es más que azúcar sintáctico para el encaje de
patrones, se define como:
\begin{align*}
  \funfont{cond} \ \texttt{true} \ v \ w &\coloneqq v \\
  \funfont{cond} \ \texttt{false} \ v \ w &\coloneqq w
\end{align*}

Podemos usar incluso más azúcar y denotar la función \funfont{cond} como un \emph{if then else}:
\begin{align*}
  \funfont{update} \ & (m: \tyfont{Map} \ \alpha) \ (k: \tyfont{String}) \ (v: \alpha): \tyfont{Map} \ \alpha \coloneqq \\
                      &\lambda k' \mapsto \LeanIfElse{k = k'}{v}{m \ k'}
\end{align*}
Que no debe ser confundido con el \texttt{if then else} del lenguaje \texttt{IMP}.
Podemos introducir también una notación para denotar la \enquote{asignación} de un mapa de
forma más familiar:
\begin{notation}[Actualización de mapa]
  Denotaremos \((\funfont{update} \ m \ k \ v)\) como \(m[k \leftarrow v]\).
\end{notation}

\begin{theorem}
  Si asignamos un valor a una variable, el mapa evalúa la clave a ese valor.
  \[
  (m[k \leftarrow v]) \ k = v
  \]
\end{theorem}

\begin{theorem}
  Si asignamos un valor a una variable, el mapa no evalúa a ese valor para
  otras variables (\(k \neq k'\)).
  \[
  (m[k \leftarrow v]) \ k' = m \ k'
  \]
\end{theorem}

Estos dos teoremas son bastante intuitivos, veamos otro teorema que también
es intuitivo pero no tan fácil de demostrar. Normalmente asignar un valor a una
variable hace que el ordenador \enquote{olvide} el valor anterior. Para la
función \funfont{update} esto no es cierto por definición, pero podemos
demostrarlo fácilmente.
\begin{theorem}
  Asignar un valor a una variable, y luego asignar otro valor a la misma
  variable hace que el mapa \enquote{olvide} el primer valor.
  \[
  m[k \leftarrow v_2][k \leftarrow v_1] = m[k \leftarrow v_1]
  \]
\end{theorem}

Todas las demostraciones se encuentran en el código, pero vamos a intentar 
explicar un poco, en términos más matemáticos, la demostración de este teorema.
\begin{proof}
  Para demostrar la igualdad de ambas funciones usamos el axioma de
  \emph{extensionalidad funcional}, con lo que demostrar el teorema equivale 
  a demostrar:
  \[
  \forall k' , (m[k \leftarrow v_2][k \leftarrow v_1]) \ k' = (m[k \leftarrow v_1]) \ k'
  \]
  Ahora podemos distinguir dos casos:
  \begin{itemize}
    \item[(\(k = k'\))] 
      En este caso la igualdad se cumple al \enquote{desenrollar} la
      definición de la función \funfont{update}:
      \[
        \LeanIfElse{k = k'}{v_1}{\LeanIfElse{k = k'}{v_2}{m \ k'}} =
        \LeanIfElse{k = k'}{v_1}{m \ k'}
      \]
      notando que \(k = k'\) es cierto, podemos simplificar la expresión a \(v_1 = v_1\).
    \item[(\(k \neq k'\))] Este caso es análogo.
  \end{itemize}
\end{proof}

Las demostraciones de los siguientes teoremas se realizan de forma parecida,
sin más que usar la extensionalidad funcional y el encaje de patrones.

\begin{theorem}
  El orden de las asignaciones no importa, si asignamos un valor a una
  variable y luego asignamos otro valor a otra variable (\(k_1 \neq k_2\)), el
  resultado es el mismo.
  \[
  m[k_2 \leftarrow v_2][k_1 \leftarrow v_1] = m[k_1 \leftarrow v_1][k_2 \leftarrow v_2]
  \]
\end{theorem}

\begin{theorem}
  Si asignamos a una variable el valor que ya tiene, el mapa no cambia.
  \[
  m[k \leftarrow m \ k] = m
  \]
\end{theorem}

Con estas propiedades demostradas podemos definir nuestro tipo \(\tyfont{State}\).
\begin{definition}[Estado]
  Un estado es un mapa que mapea variables a valores, podemos definirlo como:
  \[
  \tyfont{State} \coloneqq \tyfont{Map} \ \mathbb{Z}
  \]
\end{definition}

Pero esto nos plantea un problema, ¿qué valor tienen las variables cuando el
programa empieza? En algunos lenguajes de programación populares, como C,
cuando se leen variables antes de ser asignadas tienen un valor aleatorio, este
valor depende de lo que haya en la memoria, por lo que es \enquote{aleatorio}.
Podríamos modelar esto definiendo un estado \emph{indefinido}, pero el tipo
\(\tyfont{Map}\) que hemos definido es \emph{total}, es decir, está definido
para todas las claves. Podríamos definir un mapa parcial, que es un mapa que
no está definido para todas las claves.
\begin{definition}[Mapa parcial]
  Un mapa parcial es un mapa que no está definido para todas las claves, lo
  definimos como:
  \[
  \tyfont{PMap} \ (\alpha : \tyfont{Type}) \coloneqq \tyfont{String} \to \tyfont{Option} \ \alpha
  \]
\end{definition}

Pero esto complicaría las definiciones y las demostraciones innecesariamente, así que
usaremos la aproximación más sencilla de definir el estado inicial como un
mapa que asigna todas las variables a cero. Por supuesto, asignarles cero es
solo una convención, podríamos asignarles cualquier valor.

\begin{definition}[Estado inicial]
  El estado inicial de un programa \(\texttt{IMP}\) es la función constante que
  mapea todas las variables a cero:
  \[
  s_0 \coloneqq \lambda x \mapsto 0
  \]
\end{definition}

De esta forma evaluar el \enquote{contenido} de una \enquote{variable} es tan
sencillo como aplicar el estado a una cadena:
\[
(s_0[k \leftarrow 3][k \leftarrow 4][k \leftarrow 7]) \ k = 7
\]

\subsection{Evaluación de expresiones}

Una vez descrita la \emph{sintaxis} del lenguaje, podemos pasar a definir la
\emph{evaluación} de expresiones. Como nuestras expresiones contienen
variables, no podemos darles significado sin un estado asociado, en efecto,
consideramos la expresión \(2 + 2\), claramente su valor es \(4\), pero si
tenemos la expresión \(x + 2\) no podemos saber su valor sin conocer el
valor de \(x\). Por ello, la evaluación de expresiones se define sobre una
\emph{configuración}, que es un par de una expresión y un estado:
\(\tyfont{Aexp} \times \tyfont{State}\).

Nuestra evaluación está \emph{dirigida por la sintaxis}, es decir, para cada
\emph{constructo sintáctico} del lenguaje definiremos una regla que nos dice
como evaluarlo. Esta regla será una \emph{relación binaria} entre configuraciones
y valores, la llamaremos \(\funfont{bigStep}\) o \emph{evaluación de paso largo}.

\begin{definition}[Paso largo aritmético]
  La relación de paso largo es una relación binaria entre configuraciones y
  valores, tiene tipo:
  \[
  \funfont{bigStep} : \tyfont{Aexp} \times \tyfont{State} \to \mathbb{Z} \to \Prop
  \]  
\end{definition}

Se conoce como \emph{semántica de paso largo} porque evalúa la expresión a un
valor \enquote{en un solo paso}, desde la configuración directamente al valor.
Introducimos una notación más usada en textos clásicos de semántica formal.
\begin{notation}
  Denotamos la relación de paso largo \(\funfont{bigStep} \ (a, s) \ n\) como:
  \[
  (a, s) \bigstep n
  \]
\end{notation}

Los constructores de la relación son las \emph{reglas de evaluación}, por ejemplo el
constructor para una expresión constante es:
\[
  \texttt{val} : (n, \placeholder) \bigstep n
\]

Este constructor relaciona la configuración \((n, \placeholder)\) con el valor entero \(n\), 
donde el \(n\) de la izquierda es una expresión 
constante\footnote{En realidad deberíamos escribir \(\texttt{val} \ n\)
  para denotar que es el primer constructor del tipo \(\tyfont{Aexp}\), pero
  haremos un pequeño abuso de notación y lo omitiremos por brevedad.} y el 
\(\placeholder\) denota cualquier estado. Es importante no confundir el constructor
de la \emph{expresión sintáctica} con el de la \emph{relación de evaluación}, ya 
que aunque tienen el mismo nombre la distinción es clara por el contexto.

El constructor para una expresión variable es similar, pero esta vez necesitamos
\enquote{extraer} el valor de la variable del estado:
\[
  \texttt{var} : (x, s) \bigstep s \ x
\]

Nos planteamos ahora definir los constructores para las operaciones
aritméticas, como estas son operaciones \emph{binarias} necesitamos conocer
el valor de la expresión izquierda y derecha para evaluar la operación,
haremos esto exigiendo la evaluación de las expresiones izquierda y derecha en 
las premisas de la regla, o lo que es lo mismo nuestro constructor 
tomará otros dos constructores del mismo tipo. Por ejemplo el
constructor para la operación suma es:
\[
  \texttt{add} : \frac{(a_1, s) \bigstep n_1 \quad (a_2, s) \bigstep n_2}
   {(a_1 \ImpPlus a_2, s) \bigstep n_1 + n_2}
\]
Esta regla dice que si la configuración \((a_1, s)\) evalúa a \(n_1\) y la
configuración \((a_2, s)\) evalúa a \(n_2\), entonces la configuración
\((a_1 \ImpPlus a_2, s)\) evalúa a \(n_1 + n_2\).

Los constructores para la resta y la multiplicación son análogos:
\begin{align*}
  \texttt{sub} &: \frac{(a_1, s) \bigstep n_1 \quad (a_2, s) \bigstep n_2}
   {(a_1 \ImpMinus a_2, s) \bigstep n_1 - n_2} \\
  \texttt{mul} &: \frac{(a_1, s) \bigstep n_1 \quad (a_2, s) \bigstep n_2}
   {(a_1 \ImpTimes a_2, s) \bigstep n_1 * n_2}
\end{align*}

Podemos pensar que mientras aplicamos las reglas sintácticas del lenguaje
\enquote{construimos} la expresión a partir del texto, mientras que al aplicar 
las reglas de evaluación \enquote{descomponemos} la expresión hasta su valor.
Veamos un ejemplo de un \emph{árbol de evaluación} válido.

\begin{example}[Árbol de evaluación]\label{ex:arithmetic_evaluation}
  Evaluamos el valor de la expresión \((x + 5) + (9 - 7)\) en el estado \(s_0\)
  aplicando las \enquote{reglas} (constructores) de la relación de paso largo:
\[
  \texttt{add}: \frac{
    \texttt{add}: \displaystyle{\frac{
      \texttt{var}: (x, s_0) \bigstep 0
      \quad
      \texttt{val}: (5, s_0) \bigstep 5
    }{(x \ImpPlus 5, s_0) \bigstep 5}}
    \quad 
    \texttt{sub}: \displaystyle{\frac{
      \texttt{val}: (9, s_0) \bigstep 9
      \quad
      \texttt{val}: (7, s_0) \bigstep 7
    }{(9 \ImpMinus 7, s_0) \bigstep 2}}
  }
  {((x \ImpPlus 5) \ImpPlus (9 \ImpMinus 7), s_0) \bigstep 7}
\]
\end{example}

Este enfoque tiene una desventaja a nivel computacional y es que en cada
\enquote{paso} de la evaluación hemos tenido que \enquote{elegir} qué regla
aplicar. Puesto que en nuestro caso solo podemos aplicar una regla a cada 
constructo podemos definir una \emph{función de evaluación} en el que las
reglas se apliquen mediante encaje de patrones, y no mediante una relación
binaria. Esta función tomará una expresión y un estado y devolverá un valor.

\begin{definition}[Función de evaluación]
  La función de evaluación toma una expresión y un estado y devuelve un valor,
  tiene tipo:
  \[
  \funfont{eval} : \tyfont{Aexp} \to \tyfont{State} \to \mathbb{Z}
  \]
\end{definition}

Notamos que en vez de usar un producto cartesiano, podemos usar una función
\emph{currificada} (o \emph{parcialmente aplicada}) para definir la función de
evaluación. Esto es posible porque la función de evaluación es \emph{total},
es decir, está definida para todas las expresiones y estados, mediante las 
siguientes reglas:
\begin{align*}
  \funfont{eval} \ (\texttt{val} \ v) \ \placeholder &\coloneqq v \\
  \funfont{eval} \ (\texttt{var} \ x) \ s &\coloneqq s \ x \\
  \funfont{eval} \ (\texttt{add} \ a_1 \ a_2) \ s &\coloneqq \funfont{eval} \ a_1 \ s + \funfont{eval} \ a_2 \ s \\
  \funfont{eval} \ (\texttt{sub} \ a_1 \ a_2) \ s &\coloneqq \funfont{eval} \ a_1 \ s - \funfont{eval} \ a_2 \ s \\
  \funfont{eval} \ (\texttt{mul} \ a_1 \ a_2) \ s &\coloneqq \funfont{eval} \ a_1 \ s * \funfont{eval} \ a_2 \ s
\end{align*}

Introducimos notación para denotar la evaluación de una expresión en un
estado.
\begin{notation}[Evaluación aritmética]
  Denotamos la evaluación de la función \(\funfont{eval} \ a \ s\) como la aplicación de la expresión \(a\)
  a un estado \(s\):
  \[
  a \ s
  \]
\end{notation}

Podemos entonces evaluar el ejemplo~\ref{ex:arithmetic_evaluation} directamente
con la función \(\funfont{eval}\):
\[
  ((x \ImpPlus 5) \ImpPlus (9 \ImpMinus 7)) \ s_0 = 7
\]

La ventaja de este enfoque es evidente, ya que podemos calcular el resultado
directamente, en vez de solo decir si un \enquote{árbol de evaluación} es válido o no.
Todo esto suponiendo que ambas definiciones son equivalentes, afortunadamente lo son.
\begin{theorem}
  La relación de paso largo y la función \(\funfont{eval}\) son equivalentes, es decir, para todas
  las expresiones \(a\), estados \(s\) y valores \(n\), tenemos:
  \[
  ((a, s) \bigstep n) = (a \ s = n)
  \]
\end{theorem}

\begin{proof}
Aplicamos el axioma de \emph{extensionalidad proposicional} para transformar la
igualdad en una doble implicación, y luego probamos cada implicación por separado.

Comenzamos probando la implicación de izquierda a derecha
\((a, s) \bigstep n \implies a \ s = n\). Procedemos por inducción
sobre la relación de evaluación, los casos para los constructores \(\texttt{val}\)
y \(\texttt{var}\) son inmediatos. En el caso de \(\texttt{add}\) la
expresión tiene la forma \(a_1 + a_2\), con lo que tenemos que probar que
\[
((a_1 \ImpPlus a_2) \ s) = n_1 + n_2
\]
para algún \(n_1\) y \(n_2\), por la hipótesis de inducción tenemos que
\(a_1 \ s = n_1\) y \(a_2 \ s = n_2\), y así solo tenemos que aplicar la
definición de la función \(\funfont{eval}\) para probar la igualdad.
Los casos para la resta y la multiplicación son análogos.

La implicación de derecha a izquierda \((a \ s = n) \implies (a, s) \bigstep n\)
se demuestra por inducción sobre la expresión aritmética, los casos para
\(\texttt{val}\) y \(\texttt{var}\) son inmediatos, así que nos enfocamos en
\(\texttt{add}\). Como antes tenemos que probar que
\((a_1 \ImpPlus a_2, s) \bigstep n\), con las hipótesis de inducción:
\[
\begin{gathered}
  \forall n, a_1 \ s = n \implies (a_1, s) \bigstep n \qquad
  \forall n, a_2 \ s = n \implies (a_2, s) \bigstep n
\end{gathered}
\]
si elegimos que \(n\) es \(a_1 \ s\) y \(a_2 \ s\) respectivamente podemos construir
el árbol de evaluación:
\[
  \texttt{add}: \frac{
    (a_1, s) \bigstep a_1 \ s
    \quad
    (a_2, s) \bigstep a_2 \ s
  }
  {
   (a_1 \ImpPlus a_2, s) \bigstep a_1 \ s + a_2 \ s
  }
\]
Con lo que basta aplicar la hipótesis \((a_1 \ImpPlus a_2) \ s = n\) para obtener el resultado. Los casos para
la resta y la multiplicación son análogos.
\end{proof}

Presentamos las reglas de evaluación de expresiones booleanas sin mucho detalle,
ya que son análogas a las expresiones aritméticas, las cuales hemos presentado
primero ya que algunas de las reglas para expresiones booleanas dependen de la
evaluación de expresiones aritméticas, (\texttt{<=}) y (\texttt{==}).
\[
\begin{gathered}
  \texttt{true}: (\texttt{true}, \placeholder) \bigstep \texttt{true} \quad
  \texttt{false}: (\texttt{false}, \placeholder) \bigstep \texttt{false} \\
  \texttt{not}: \frac{(b, s) \bigstep b'}{(\ImpNot b, s) \bigstep \lnot b'} \quad
  \texttt{and}: \frac{(b_1, s) \bigstep b_1' \quad (b_2, s) \bigstep b_2'}
    {(b_1 \ \ImpAnd \ b_2, s) \bigstep b_1' \land b_2'} \quad
  \texttt{or}: \frac{(b_1, s) \bigstep b_1' \quad (b_2, s) \bigstep b_2'}
    {(b_1 \ \ImpOr \ b_2, s) \bigstep b_1' \lor b_2'} \\
  \texttt{eq}: (a_1 \ \ImpEq \ a_2, s) \bigstep a_1 \ s = a_2 \ s \quad
  \texttt{le}: (a_1 \ \ImpLe \ a_2, s) \bigstep a_1 \ s \leq a_2 \ s
\end{gathered}
\]

Notamos que en las últimas dos reglas estamos usando la función
(\funfont{eval}) para evaluar las expresiones aritméticas, podríamos haber
usado la relación de paso largo, pero este enfoque es más sencillo, y como hemos
visto antes equivalente. Todo lo que hemos dicho sobre la función
(\funfont{eval}) y su equivalencia con la relación de paso largo se aplica a las
expresiones booleanas.

\subsection{Evaluación de comandos}

Hemos visto en la sección anterior que lo más cómodo para definir la semántica de
expresiones aritméticas y booleanas es usar la función de evaluación
(\funfont{eval}). Vamos a intentar hacer lo mismo con los comandos, las reglas
de nuestra función de evaluación serían algo como:
\begin{align*}
  \funfont{eval} \ \ImpSkip \ s &\coloneqq s \\
  \funfont{eval} \ \ImpAssign{x}{a} \ s &\coloneqq s[x \leftarrow a \ s] \\
  \funfont{eval} \ (\ImpConcat{c_1}{c_2}) \ s &\coloneqq \funfont{eval} \ c_2 \ (\funfont{eval} \ c_1 \ s) \\
  \funfont{eval} \ (\ImpIfElse{b}{c_1}{c_2}) \ s &\coloneqq
    \LeanIfElse{b \ s}{\funfont{eval} \ c_1 \ s}{\funfont{eval} \ c_2 \ s} \\
  \funfont{eval} \ (\ImpWhile{b}{c}) \ s &\coloneqq
    \LeanIfElse{b \ s}{\funfont{eval} \ (\ImpWhile{b}{c}) \ (\funfont{eval} \ c \ s)}{s}
\end{align*}
El problema de esta función es que no es en realidad una función, en el 
sentido de que no es \emph{bien fundada}. Consideramos el bucle \(\ImpWhile{\texttt{true}}{\ImpSkip}\),
claramente aplicando la última regla de la función de evaluación nunca
obtenemos un resultado, ya que la condición es siempre verdadera\footnote{
  Demostramos esto rigurosamente en~\ref{sec:non_termination}.}.
Por tanto deberemos recurrir nuevamente al paso largo, más tarde veremos como
podemos definir una \enquote{función} que evalúe los comandos, que es lo que se
conoce como \emph{semántica denotacional}.

Definimos la semántica de paso largo de los comandos como una relación, como
antes, esta vez asociando configuraciones a estados. Damos primero
la definición de los constructores más sencillos.
\begin{align*}
  \ImpSkip &: (\ImpSkip, s) \bigstep s \\
  \texttt{ass} &: (v = a, s) \bigstep s[v \leftarrow a \ s] \\
  \texttt{cat} &: \frac{
    (c_1, s_1) \bigstep s_2 \quad (c_2, s_2) \bigstep s_3
  }{
    (\ImpConcat{c_1}{c_2}, s_1) \bigstep s_3
  }
\end{align*}
En la concatenación introducimos un estado intermedio \(s_2\) que no 
aparece en el resultado final, de esta forma tenemos en cuenta los 
\enquote{pasos intermedios} en la evaluación del lado izquierdo.

Las reglas para la estructura \texttt{if} requieren dos constructores,
uno para el caso en que la condición es falsa:
\[
\texttt{ifFalse} : \frac{b \ s_1 = \texttt{false}  \quad (c_2, s_1) \bigstep s_2}
{(\ImpIfElse{b}{c_1}{c_2}, s_1) \bigstep s_2}
\]
y otro cuando la condición es cierta:
\[
\texttt{ifTrue} : \frac{b \ s_1 = \texttt{true}  \quad (c_1, s_1) \bigstep s_2}
{(\ImpIfElse{b}{c_1}{c_2}, s_1) \bigstep s_2}
\]
La regla de evaluación para \texttt{while} sigue el mismo principio pero con un enfoque más
recursivo, comenzamos con el caso en que la condición es falsa, que es el más sencillo:
\[
\texttt{whileFalse} : \frac{b \ s = \texttt{false}}
  {(\ImpWhile{b}{c}, s) \bigstep s}
\]
pero cuando la condición es cierta, tenemos que hacer una recursión desde otro árbol de evaluación
donde el \enquote{paso anterior} ya ha sido aplicado:
\[
\texttt{whileTrue} : \frac{b \ s_1 = \texttt{true} \quad (c, s_1) \bigstep s_2 \quad
  (\ImpWhile{b}{c}, s_2) \bigstep s_3}
{(\ImpWhile{b}{c}, s_1) \bigstep s_3}
\]
Mientras que las reglas de paso largo para expresiones pueden ser aplicadas
\enquote{automáticamente}, las reglas de paso largo para comandos requieren un enfoque
más \enquote{manual}, ya que tenemos que \enquote{elegir} el constructor correcto
en el flujo de control del programa.

Con los abusos de notación habituales, podemos comprobar la evaluación de un
programa como este:
\begin{example}[Árbol de evaluación de un programa]
  Evaluamos un programa sencillo que contiene una bifurcación en el estado \(s_0\):
\begin{minted}{rust}
  x = 2;
  if x <= 1 {
    y = 3
  } else {
    z = 4
  }
\end{minted}
  Para aligerar la notación denotaremos el programa como \(P\), y la estructura 
  \texttt{if} como \(\text{Ifelse}\), entonces:
  \[
  \texttt{cat}: \frac{
    \texttt{ass}: (\ImpAssign{x}{2}, s_0) \bigstep s_0[x \leftarrow 2]
    \quad
    \texttt{ifFalse}: \displaystyle{\frac{
      \texttt{ass}: (\ImpAssign{z}{4} , s_0[x \leftarrow 2]) \bigstep s_0[x \leftarrow 2][z \leftarrow 4]
    }
    {
      (\text{Ifelse}, s_0[x \leftarrow 2]) \bigstep s_0[x \leftarrow 2][z \leftarrow 4]
    }}
  }{
    (P, s_0) \bigstep s_0[x \leftarrow 2][z \leftarrow 4]
  }
  \]
\end{example}
Como podemos observar la \enquote{demostración} de que el programa termina
en un estado dado se sigue directamente del árbol de evaluación creado por los
constructores de la semántica de paso largo.

\subsection{Reglas de reescritura}

Presentamos ahora algunas \enquote{reglas de reescritura} para los constructos del
lenguaje que nos permiten simplificar las demostraciones. Solamente enunciaremos
los teoremas, ya que las demostraciones se pueden encontrar en el código y no
tienen mucho interés.
\begin{theorem}
  Usar el comando \ImpSkip{} (no hacer nada) no cambia el estado:
  \[
  ((\ImpSkip, s_1) \bigstep s_2) = (s_1 = s_2)
  \]
\end{theorem}

\begin{theorem}
  Concatenar dos programas es equivalente a ejecutar el primero y luego el segundo:
  \[
  ((\ImpConcat{c_1}{c_2}, s_1) \bigstep s_3) = 
  \exists s_2, \left((c_1, s_1) \bigstep s_2 \land (c_2, s_2) \bigstep s_3\right)
  \]
\end{theorem}

Las reglas de reescritura para el constructo \texttt{if} pueden definirse de
dos formas equivalentes, qué regla aplicar depende del contexto, ya que
cualquiera de las dos podría ser más fácil de usar dependiendo de la situación.
\begin{theorem}\label{thm:if_rewrite_1}
  El constructo \texttt{if} puede reescribirse como:
  \[
  ((\ImpIfElse{b}{c_1}{c_2}, s_1) \bigstep s_2) = 
  \LeanIfElse{b \ s_1}{(c_1, s_1) \bigstep s_2}{(c_2, s_1) \bigstep s_2}
  \]
\end{theorem}
O lo que es lo mismo:
\begin{theorem}\label{thm:if_rewrite_2}
  El constructo \texttt{if} puede reescribirse como:
  \[
  ((\ImpIfElse{b}{c_1}{c_2}, s_1) \bigstep s_2) =
  (\LeanIfElse{b \ s_1}{c_1}{c_2}, s_1) \bigstep s_2
  \]
\end{theorem}

En los teoremas anteriores tenemos que tener cuidado de no confundir el
\enquote{if} del lenguaje \texttt{IMP} y el del meta-lenguaje (en este caso Lean).
Para ayudar en esta distinción usaremos una fuente \texttt{mono espaciada} para los
constructos del lenguaje, y una fuente normal para los constructos del
meta-lenguaje.

\begin{theorem}\label{thm:while_rewrite}
  El constructo \texttt{while} puede reescribirse como:
  \begin{align*}
  &((\ImpWhile{b}{c}, s_1) \bigstep s_3) = \\
  &\quad \text{if} \ b \ s_1 \\
  &\quad \text{then} \ \exists s_2, ((c, s_1) \bigstep s_2) \land ((\ImpWhile{b}{c}, s_2) \bigstep s_3) \\
  &\quad \text{else} \ s_1 = s_3
  \end{align*}
\end{theorem}

\subsection{Equivalencia de comportamiento}

Las reglas de reescritura anteriores son útiles en las demostraciones ya que
podemos sustituir un programa por otro que es \enquote{exactamente el mismo}, pero
podemos también definir una noción de igualdad más laxa que nos permita decir que
dos programas sintácticamente diferentes son \enquote{el mismo}. Esto es
especialmente útil cuando queremos demostrar que ciertas optimizaciones no
cambian el comportamiento de un programa.

\begin{definition}[Equivalencia de comportamiento]
  Dos programas se comportan igual si para todos los estados iniciales:
  \[
  c_1 \sim c_2 \coloneqq \forall s_1 \ s_2, (c_1, s_1) \bigstep s_2 \iff (c_2, s_1) \bigstep s_2
  \]
\end{definition}

Es fácil ver que esta relación es una relación de equivalencia, es decir, es
reflexiva, simétrica y transitiva. Con esto podemos, por ejemplo, demostrar que
agregar un número arbitrario de \ImpSkip{} a un programa no cambia su
comportamiento.
\begin{theorem}
  Agregar \ImpSkip{} a la izquierda de un programa no cambia su comportamiento:
  \[
  \ImpConcat{\ImpSkip}{c} \sim c
  \]
\end{theorem}

Ahora podemos \enquote{optimizar} el \ImpSkip{} de un programa, pero solo
a la izquierda, para eliminarlo completamente sin consecuencias necesitamos el
siguiente teorema.
\begin{theorem}
  Agregar \ImpSkip{} a la derecha de un programa no cambia su comportamiento:
  \[
  \ImpConcat{c}{\ImpSkip} \sim c
  \]
\end{theorem}

Una equivalencia más interesante que refleja el teorema~\ref{thm:while_rewrite} es
el \emph{desenrollado de bucles}, una optimización común en compiladores, usada
cuando la memoria es más barata que la computación.
\begin{theorem}
  El comando \texttt{while} puede desenrollarse:
  \[
  \ImpWhile{b}{c} \sim 
  \ImpIfElse{b}{\ImpConcat{c}{\ImpWhile{b}{c}}}{\ImpSkip}
  \]
\end{theorem}

\begin{proof}
  Por la definición de equivalencia el problema se reduce a:
  \[
    \forall s_1 \ s_3, (\ImpWhile{b}{c}, s_1) \bigstep s_3 \iff
    (\ImpIfElse{b}{\ImpConcat{c}{\ImpWhile{b}{c}}}{\ImpSkip}, s_1) \bigstep s_3
  \]
  Podemos aplicar el teorema~\ref{thm:if_rewrite_2} al lado derecho de la
  equivalencia:
  \[
    (\ImpWhile{b}{c}, s_1) \bigstep s_3 \iff
    (\LeanIfElse{b \ s_1}{\ImpConcat{c}{\ImpWhile{b}{c}}}{\ImpSkip}, s_1) \bigstep s_3
  \]
  Y ahora basta con probar ambas implicaciones por separado.
  \begin{itemize}
    \item[(\(\implies\))] 
    Comprobamos ambos casos para el constructor \texttt{while} en el lado izquierdo de la implicación.
    \begin{itemize}
      \item[(\(b \ s_1\))]
      Si la condición es cierta, el problema se simplifica a:
      \[
        \ImpConcat{c}{\ImpWhile{b}{c}}, s_1 \bigstep s_3
      \]
      y como sabemos que la precondición es cierta podemos aplicar el constructor
      \(\texttt{whileTrue}\) para obtener el árbol de derivación del lado izquierdo de la equivalencia:
      \[
        \texttt{whileTrue}: \frac{
          b \ s_1 = \texttt{true}
          \quad
          (c, s_1) \bigstep s_2 
          \quad 
          (\ImpWhile{b}{c}, s_2) \bigstep s_3
        }{
          (\ImpWhile{b}{c}, s_1) \bigstep s_3
        }
      \]
      Basta con aplicar la nueva hipótesis de \(\texttt{whileTrue}\) para
      obtener el resultado:
      \[
        \texttt{cat}: \frac{
          (c, s_1) \bigstep s_2
          \quad
          (\ImpWhile{b}{c}, s_2) \bigstep s_3
        }{
          (\ImpConcat{c}{\ImpWhile{b}{c}}, s_1) \bigstep s_3
        }
      \]
      \item[(\(\lnot b \ s_1\))]
      El problema se reduce a:
      \[
        (\ImpSkip, s_1) \bigstep s_3
      \]
      procedemos como antes y aplicamos el constructor \(\texttt{whileFalse}\):
      \[
        \texttt{whileFalse}: \frac{
          b \ s_1 = \texttt{false}          
        }{
          (\ImpWhile{b}{c}, s_1) \bigstep s_1
        }
      \]
      con lo que tenemos que \(s_1 = s_3\) y hemos terminado.
    \end{itemize}

    \item[(\(\impliedby\))]
    Esta dirección se demuestra también por casos sobre la condición
    del comando \texttt{if}, pero es aún más sencilla que antes, ya que podemos
    simplificar la precondición directamente.
  \end{itemize}
\end{proof}

\subsection{No terminación}\label{sec:non_termination}

Como dijimos antes cuando intentamos definir una función \funfont{eval} para
evaluar comandos, existen programas de \texttt{IMP} que no terminan, o lo que
es lo mismo en términos de la relación de paso largo, no existe ninguna 
relación entre el programa en el estado dado y otro estado.

\begin{definition}[No terminación]
  Decimos que un programa no termina en un estado \(s\) si no existe ningún
  estado \(s'\) tal que el programa pueda evaluarse a \(s'\), es decir:
  \[
   \forall s', (c, s) \centernot \bigstep s'
  \]
\end{definition}

Ahora demostramos que existen programas en \texttt{IMP} que no terminan, para
ello elegimos de hecho un programa que no termina sin importar el estado
en el que se ejecute.

\begin{theorem}\label{thm:non_termination}
  El programa \(\ImpWhile{b}{c}\), donde \(b \sim \texttt{true}\) no termina,
  independientemente del estado inicial:
  \[
  \forall s_1 \ s_2, (\ImpWhile{b}{c}, s_1) \centernot \bigstep s_2
  \]
\end{theorem}

\begin{proof}
  En realidad el \(\centernot \bigstep\) es azúcar sintáctico para
  la negación de la relación de paso largo, así que el teorema es en realidad:
  \[
  (\ImpWhile{b}{c}, s_1) \bigstep s_2 \implies \bot
  \]
  Ahora tenemos que usar un pequeño truco, primero generalizamos la configuración
  a la izquierda, llamándola \text{hconf}:
  \[\text{hconf} \coloneqq (\ImpWhile{b}{c}, s_1)\] 
  podemos escribir el teorema como:
  \[
  \text{hconf} \bigstep s_2 \implies \bot
  \]
  Ahora aplicamos el principio de inducción sobre la configuración, es decir,
  demostramos el teorema para todas las formas posibles de \text{hconf}.
  Los únicos casos que nos interesan son los de \texttt{whileTrue} y
  \texttt{whileFalse}, ya que para el resto de los constructores la
  igualdad es falsa por definición, véase \(\text{hconf} = (\ImpSkip, s_1)\).

  Para el caso \texttt{whileTrue} tenemos que probar que:
  \[
    (\ImpWhile{b}{c}, s_1) = (\ImpWhile{b}{c}, s_1')
  \]
  la hipótesis de inducción toma la forma:
  \[
    \forall x, (\ImpWhile{b}{c}, x) = (\ImpWhile{b'}{c'}, s_2')
  \]
  así que basta elegir \(x = s_2\) para obtener el resultado. 
  
  Para el caso \texttt{whileFalse} no necesitamos usar inducción, de hecho si
  \text{hconf} toma esta forma significa que la condición del \texttt{while} es
  falsa pero \(b \sim \texttt{true}\) por hipótesis así que tenemos una
  contradicción.
\end{proof}

\subsection{Determinismo}

Otra propiedad importante de \texttt{IMP} que podemos probar usando el paso largo
es que la evaluación de un programa es determinista, es decir, dado un
programa y un estado solo hay un estado al que el programa puede evaluarse.

\begin{definition}[Determinismo]\label{thm:deterministic}
  Diremos que un lenguaje es \emph{determinista} si se cumple:
  \[
  \forall c \ s \ s_1 \ s_2, ((c, s) \bigstep s_1 \land (c, s) \bigstep s_2) \implies s_1 = s_2
  \]
\end{definition}

\begin{theorem}
  El lenguaje \texttt{IMP} es determinista.
\end{theorem}

\begin{proof}
  Haremos la demostración por inducción sobre la relación de paso largo
  \((c, s) \bigstep s_1\), los casos para \ImpSkip{} y la asignación son
  inmediatos, así que pasamos al caso de la concatenación. Necesitamos 
  la hipótesis de inducción:
  \begin{align*}
    \forall x, (c_1', s_1') &\bigstep x \implies s_3 = x \\
    \forall y, (c_2', s_2') &\bigstep y \implies s_1 = y
  \end{align*}
  El lado derecho de la conjunción toma la forma
  \((\ImpConcat{c_1'}{c_2'}, s_1') \bigstep s_2\) así que tenemos que tener
  sus precondiciones, es decir:
  \[
    \texttt{ass}: \frac{
      (c_1', s_1') \bigstep s_2' 
      \quad 
      (c_2', s_2') \bigstep s_2
    }{
      (\ImpConcat{c_1'}{c_2'}, s_1') \bigstep s_2
    }
  \]
  Ahora aplicamos la primera hipótesis de inducción en la precondición
  izquierda, obteniendo \(s_3 = s_2'\), con lo que la segunda precondición
  se transforma en \((c_2', s_3) \bigstep s_2\) y aplicando la
  segunda hipótesis de inducción obtenemos el resultado \(s_1 = s_2\).

  El resto de los casos son análogos.
\end{proof}

\section{Semántica estructural}

Otra forma de definir la semántica del lenguaje es la semántica de \emph{estructural}
o \emph{de paso corto}. En este caso en lugar de obtener directamente el valor
de la expresión tomamos un \enquote{único paso} que relaciona dos configuraciones,
a la izquierda tenemos el estado antes de tomar el paso y a la derecha el
estado después de tomarlo, al hacer esto podemos observar los \emph{estados inter
medios} de la ejecución de un programa, así que cada paso de esta semántica es
\enquote{atómico}. Esto es muy útil cuando queremos probar propiedades de
programas paralelos o concurrentes, aunque \texttt{IMP} no lo sea.
\begin{definition}[Semantica estructural]
  La semántica estructural de un programa es una relación binaria 
  \(\text{SmallStep} : \tyfont{Com} \times \tyfont{State} \to \tyfont{Com} \times \tyfont{State} \to \Prop\),
  que denotamos como:
  \[
    (c_1, s_1) \rightsquigarrow (c_2, s_2) \coloneq \text{SmallStep} \ (c_1, s_1) \ (c_2, s_2)
  \]
\end{definition}

En el caso de \texttt{IMP} los constructores de la relación son:
\[
\begin{gathered}
  \texttt{ass}: (\ImpAssign{v}{a}, s) \rightsquigarrow (\ImpSkip, s[v \leftarrow s]) \quad
  \texttt{skipCat}: (\ImpConcat{\ImpSkip}{c}, s) \rightsquigarrow (c, s) \quad
  \texttt{cat}: \frac{(c, s) \rightsquigarrow (c', s')}{(\ImpConcat{c}{c''}, s) \rightsquigarrow (c', c'', s')} \\
  \texttt{ifElse}: (\ImpIfElse{b}{c}{c'}, s)
    \rightsquigarrow (\LeanIfElse{b \ s}{c}{c'}, s) \\
  \texttt{whileLoop}: (\ImpWhile{b}{c}, s)
    \rightsquigarrow (\LeanIfElse{b \ s}{\ImpConcat{c}{\ImpWhile{b}{c}}}{\ImpSkip}, s) 
\end{gathered}
\]
Estas reglas capturan la \enquote{ejecución} de los comandos de izquierda a derecha
en una forma secuencial, lo que simula más de cerca la ejecución de un programa
en un computador. Por supuesto, podríamos haber definido las reglas del 
\enquote{if} y \enquote{while} de otras formas dependiendo de la \enquote{granularidad}
que queramos en la semántica, evaluar el boolean podría ser un estado y seleccionar la rama
otro, por ejemplo.

Veamos un ejemplo de como funciona la semántica estructural, para ello
evaluamos un programa sencillo.
\begin{example}\label{ex:small_step}
  Consideramos el programa \texttt{x = 1; while x <= 1 \{x = x + 1\}},
  para aligerar la notación denotaremos el bloque \enquote{while} como \(W\).
  El árbol de evaluación de este programa es:
  \[
  \texttt{cat}: \frac{
    \texttt{ass}: (\ImpAssign{x}{1} , s_0) \rightsquigarrow (\ImpSkip, s_0[x \leftarrow 1])
  }
  {(\ImpConcat{\ImpAssign{x}{1}}{W}, s_0)
  \rightsquigarrow
  (\ImpConcat{\ImpSkip}{W}, s_0[x \leftarrow 1])
  }
  \]
\end{example}

In this case the computation ends when we can't apply any rule to the program, 
that is when we have a configuration of the form \((\ImpSkip, \placeholder)\).
We will define the execution of a program from beginning to end with the small-step
semantics by using its \emph{reflexive transitive closure}.

\subsection{Reflexive transitive closure}

The reflexive transitive closure takes a binary relation and gives us he 
\enquote{minimal extension} of this relation that is \emph{reflexive}
and \emph{transitive}. Formally we can define it as a function that 
takes a binary relationship and return another.

\begin{definition}[Reflexive transitive closure]
  The reflexive transitive closure of a binary relation \(\mathcal{R}\) is the smallest 
  relation that contains \(\mathcal{R}\) and is both reflexive and transitive. 
  We can construct a relation \(\mathcal{R}^*\) from \(\mathcal{R}\) as an inductive type
  \(\tyfont{ReflTrans}: (\alpha \to \alpha \to \Prop) \to
  (\alpha \to \alpha \to \Prop)\) and denote it as 
  \(\mathcal{R}^* \coloneqq \tyfont{ReflTrans} \ \mathcal{R}\).

  The type \(\tyfont{ReflTrans}\) has two constructors:
  \[
    \texttt{refl}: a \mathcal{R}^* a
    \qquad 
    \texttt{tail}: \frac{
      a \mathcal{R}^* b
      \quad b \mathcal{R} c
    }{a \mathcal{R}^* c}
  \]
\end{definition}

As we can see this relation is not transitive by definition, but it is easy
to prove by induction.

\begin{theorem}
  The reflexive transitive closure is transitive, that is given 
  \(x \mathcal{R}^* y\) and \(y \mathcal{R}^* z\) we have \(x \mathcal{R}^* z\).
\end{theorem}

Another immediate but important property is that the relation itself contained
in the reflexive transitive closure.

\begin{theorem}
  The relation is contained in its reflexive transitive closure, that is
  \(a \mathcal{R} b \implies x \mathcal{R}^* y\).
\end{theorem}

Of course, we can not only append to the \enquote{tail} of the relation, but also
push to the \enquote{head}.

\begin{theorem}
  We can push a relation to the head of its reflexive transitive closure, that is,
  if we have \(a \mathcal{R} b\) and \(b mathcal{R}^* c\) then
  \(a \mathcal{R}^* c\).
\end{theorem}

This will become important soon, as we can define a principle of induction
on the \emph{head} of the relation, note that the structural induction
on the constructors works only on the \emph{tail} of the relation.

\begin{theorem}
  We can perform induction on the \emph{head} of the relationship, with the base
  case being the reflexivity, and the inductive step being the transitivity on
  the \emph{head} of the relation.
\end{theorem}

Another important theorem that we will use is the \emph{lifting} of a reflexive
transitive closure via a function, that is, if we have a relation between two
elements of a set we can lift this relation to the image of a function.

\begin{theorem}\label{thm:rel_lift}
  We can lift a reflexive transitive \(x \mathcal{R}^* y\) closure through a function
  \(f: \alpha \to \beta\). 
  Suppose that one relationship implies the other: 
  \[h: \forall a \ b, a \mathcal{R} b \to (f \ a) \mathcal{P} (f \ b),\]
  then we have \((f \ a) \mathcal{P}^* (f \ b)\). We can represent this as the 
  diagram~\ref{fig:rel_lift}, where
  the hypothesis \(h\) is expressed by the fact that the diagram commutes.
\end{theorem}

\begin{figure}[htb]\label{fig:rel_lift}
  \centering
  \begin{tikzcd}
    \alpha \times \alpha \arrow[dd, "f \times f"] \arrow[dr, "\mathcal{R}^*"] & \\
     & \Prop \\
    \beta \times \beta \arrow[ur, "\mathcal{P}^*"] &
  \end{tikzcd}
  \caption{Lifting of a reflexive transitive closure through a function.
           In reality, instead of \(\alpha \times \alpha\)
           and \(\beta \times \beta\) we should use
           the \emph{exponential objects} \(\alpha^\alpha\) and 
            \(\beta^\beta\). 
  }
\end{figure}

\begin{proof}
  We proceed by induction on \(a \mathcal{R}^* b\), the base case 
  \(\texttt{refl}: a \mathcal{R}^* a\) implies that \(a = b\), so
  \((f \ a) \mathcal{P}^* (f \ a)\)
  which is immediate by the reflexivity of the relation. For the inductive step:
  \[
    \texttt{tail}: \frac{
      a \mathcal{R}^* b 
      \quad 
      b \mathcal{R} c
    }{a \mathcal{R}^* c}
  \]
  we have to prove that \((f \ a) \mathcal{P}^* (f \ c)\), applying the
  constructor:
  \[
    \funfont{tail}: \frac{
      \text{ih}: (f \ a) \mathcal{P}^* (f \ b)
      \quad
      h: \displaystyle{\frac{
        b \mathcal{R} c
      }{
        (f \ b) \mathcal{P} (f \ c)
      }}
    }{
      (f \ a) \mathcal{P}^* (f \ c)
    }
  \]
  where \(\text{ih}\) is the induction hypothesis, we get the result.
\end{proof}

Now we can define the reflexive transitive closure of the small-step semantics.

\begin{notation}
  We will denote the reflexive transitive closure of the small-step semantics as
  \(\rightsquigarrow^*\).
\end{notation}

Now we can revisit the example~\ref{ex:small_step} with the reflexive transitive
closure of the small-step semantics.

\begin{example}
  As before, we consider the program \texttt{x = 1; while x <= 1 \{x = x + 1\}},
  to ease reading we will denote the \enquote{while} block as \(W\). We have the 
  partial derivation tree:
  \[
    \funfont{head}: \frac{
      \funfont{cat}: \displaystyle{\frac{
        \funfont{ass}: (\texttt{x = 1} , s_0) \rightsquigarrow (\ImpSkip, s_0[x \leftarrow 1])
      }
      {(\ImpConcat{\ImpAssign{x}{1}}{W}, s_0)
      \rightsquigarrow
      (\ImpConcat{\ImpSkip}{W}, s_0[x \leftarrow 1])
      }}
      \quad
      \funfont{head}: \displaystyle{\frac{
        \funfont{skipcat}: \ImpSkip
        \quad
        \vdots
      }{
        (W, s_0[x \leftarrow 1]) \rightsquigarrow^* (W, s_0[x \leftarrow 1][x \leftarrow 2])
      }}
    }{
      (\ImpConcat{\ImpAssign{x}{1}}{W}, s_0) \rightsquigarrow^* (\ImpConcat{\ImpSkip}{W}, s_0[x \leftarrow 1][x \leftarrow 2])
    }
  \]
  By successive applications of the constructors we get the result.
\end{example}

We preset now some theorems that will be very useful when proving the equivalence of
the big and small-step semantics, paying attention to the constructors of the small-step
relation we can see that the only \enquote{special} one is \texttt{cat}, since
it is the only one that requires an argument. For this reason all the following
theorems will give us analogous in \(\rightsquigarrow^*\) to the 
big-step semantics.

\begin{theorem}
  If an expression evaluates to a state concatenating that expression with another
  has that state as an intermediate. That is if \((c, s) \rightsquigarrow^* (\ImpSkip, s')\)
  then \((\ImpConcat{c}{c'}, s) \rightsquigarrow^* (\ImpConcat{\ImpSkip}{c'}, s')\).
\end{theorem}

\begin{proof}
  The proof follows immediately by Theorem~\ref{thm:rel_lift}.
  The relations \(\mathcal{R}\) and \(\mathcal{P}\) are both the small-step semantics
  \(\rightsquigarrow\), and the function \(f\) is the concatenation of commands:
  \(f: \texttt{Com} \times \texttt{State} \to \texttt{Com} \times \texttt{State}\)
  with \(f \coloneqq \lambda (c, s) \mapsto (\ImpConcat{c}{c'}, s)\). Our hypothesis \(h\)
  is the concatenation of the small-step semantics, that is:
  \[h: \Forall_{(c, s), (d, t)} (c, s) \rightsquigarrow (d, t) \to (\ImpConcat{c}{c'}, s) \rightsquigarrow (\ImpConcat{d}{c'}, t)
     \coloneqq \lambda x \mapsto \funfont{cat} \ x\]
  and finally \((c, s) \rightsquigarrow^* (\ImpSkip, s')\) by hypothesis.
\end{proof}

\begin{theorem}
  Concatenating an expression that evaluates to an intermediate configuration to
  another that evaluates in that configuration gives us the configuration of the
  second expression. That is if \((c, s) \rightsquigarrow^* (\ImpSkip, s')\)
  and \((c', s') \rightsquigarrow^* (\ImpSkip, s'')\) then
  \((\ImpConcat{c}{c'}, s) \rightsquigarrow^* (\ImpSkip, s'')\).
\end{theorem}

\begin{theorem}
  Concatenating an evaluation that doesn't change the state with another
  evaluation is equivalent to evaluating the second expression. That is
  if \((c, s) \rightsquigarrow^* (\ImpSkip, s)\) then
  \((\ImpConcat{c}{c'}, s) \rightsquigarrow^* (c', s)\).
\end{theorem}

\section{Relational denotational semantics}

We follow the construction presented  in~\cite[Chapter~11]{semantics_with_isabelle} 
and define a \emph{relational} denotational semantics for the language \texttt{IMP}.

As we saw before, defining a function that evaluates a program is not trivial,
in fact this function can never be \emph{total} as there are programs that don't
terminate (Theorem~\ref{thm:non_termination}). The function then
has to necessarily be \emph{partial}, that is, it doesn't return a value for
all inputs. This is a problem as in type theory all our functions are
total. To circumvent this problem we could use the \texttt{Option} monad,
and return \texttt{none} when the program doesn't terminate:
\[
  \funfont{eval} : \tyfont{Com} \to \tyfont{State} \to \tyfont{Option} \ \tyfont{State}.
\]

Instead of this, we will use the classic definition of relations as sets
\(\tyfont{Set} \ (\alpha \times \beta)\), this will allow us to use the 
partial ordering of subsets to define the semantics of the language via
the \emph{Knaster-Tarski} fixpoint theorem. Of course, relations are not 
partial functions, to circumvent this we will introduce the concept
of \(\omega\)-chains.

\subsection{Knaster-Tarski fixpoint theorem}

The \emph{Knaster-Tarski} fixpoint theorem is a generalization of the
\emph{Banach} fixpoint theorem, it states that every monotone function
on a complete lattice has a least fixpoint. This will be essential when defining
the denotational semantics if the \texttt{while} construct. Before
we can enunciate the theorem we have to introduce some definitions.

\subsubsection{Order theory}

We begin by defining a \emph{preorder}, which is the \enquote{weakest} form of
order that we can define.

\begin{definition}[Preorder]
  A preorder is a binary relation (usually denoted by \(\leq\)) that satisfies
  the following properties:
  \begin{itemize}
    \item Reflexivity: \(a \leq a\).
    \item Transitivity: \(a \leq b \to b \leq c \to a \leq c\).
    \item Less-than iff less-than-or-equal: \(a < b \iff a \leq b \land b \not\leq a\).
  \end{itemize}
\end{definition}

Adding antisimmetry to the preorder we obtain a \emph{partial order}.

\begin{definition}[Partial order]
  A partial order is a preorder that is antisymmetric, that is, it satisfies:
  \(a \leq b \to b \leq a \to a = b\).
\end{definition}

Now we consider a set of elements with a partial order, if every set has an
an \emph{infimum} we call it a \emph{complete lattice}. Usually a complete lattice 
is defined as a partially ordered type in which every set has a \emph{least upper bound}
and a \emph{greatest lower bound}, but these definitions are interchangeable.
Although we won't prove it, the existence of one implies the existence of the other.

\begin{definition}
  A complete lattice is a partially ordered type that has an infimum in every set. 
  That is, exists a function \(\funfont{Inf}: \tyfont{Set} \ \alpha \to \alpha\)
  that returns an element of the set satisfying the following properties:
  \begin{itemize}
    \item \(\forall S, \forall x \in S, \funfont{Inf} \ S \leq x\).
    \item \(\forall S, (\forall y \in S, x \leq y) \implies x \leq \funfont{Inf} \ S\).
  \end{itemize}
\end{definition}

Additionally, this infimum is \emph{unique}.

\begin{theorem}
  The infimum of a set in a complete lattice is unique.
  That is if \(\funfont{Inf} \ s = a\) and \(\funfont{Inf} \ s = b\) then \(a = b\).
\end{theorem}

Now we need to recall the definition of subsets in~\ref{sec:sets},
as sets, with the inclusion relation, form a partial order.

\begin{theorem}
  Sets with the inclusion relation form a partial order. Where
  \(A \leq B \coloneqq A \subseteq B\), \(A < B \coloneqq A \subset B\).
  and the infimum of a set is the intersection of all the sets in the set,
  that is \(\funfont{Inf} \ S = \bigcap_{A \in S} A\).
\end{theorem}

The intersection of a set of sets is the set of elements that are in all the sets
of the set, that is:
\[
  \bigcap_{A \in S} A = \{x \mid \forall A \in S, x \in A\}
\]
Since the proofs are
straightforward we will not present them, but they can be found in the code.

\subsubsection{Order homomorphims}

Of special interest to us are the functions that preserve the order of the elements
of a lattice, these are called \emph{order homomorphisms}, and as we will see
later they admit a \emph{least fixpoint}.

\begin{definition}[Monotonic function]
  A function \(f: \alpha \to \alpha\) is monotonic if for all \(a, b: \alpha\)
  we have that \(a \leq b \implies f \ a \leq f \ b\). We call these functions
  \emph{Order homomorphisms}.
\end{definition}

\begin{notation}
  We will denote an order homomorphism as \(f: \alpha \to_o \alpha\).
\end{notation}

Some order homomorphisms are the identity function and constant functions.

\begin{theorem}
  The identity function is an order homomorphism.
\end{theorem}

\begin{theorem}
  Any constant function is an order homomorphism.
\end{theorem}

\subsubsection{Fix-points}

As we said before, monotonic functions on a complete lattice have a least fixpoint,
we will define the concept of fixpoint and least fixpoint, but first we need
to define the concept of pre-fixpoint.

\begin{definition}[Pre-fixpoint]
  A pre-fixpoint of a function \(f: \alpha \to_o \alpha\) is a point \(a: \alpha\) 
  such that \(f \ a \leq a\).
\end{definition}

Now we can define the \emph{least pre-fixpoint} of an order homomorphism.

\begin{definition}[Least pre-fixpoint]
  The least pre-fixpoint of a function \(f: \alpha \to_o \alpha\) is the infimum of 
  the set of all pre-fixpoints of \(f\).
  \[
    \funfont{lfp} \ f = \funfont{Inf} \ \{a \mid f \ a \leq a\}
  \]
\end{definition}

In the case of sets this turns into 
\(\funfont{lfp} \ f = \bigcap \{a \mid f \ a \subseteq a\}\), 
where \(f: \tyfont{Set} \ \alpha \to \tyfont{Set} \ \alpha\).

Now we present two important pre-fixpoint inequalities.

\begin{theorem}
  The least pre-fixpoint of a function is less than or equal to any
  pre-fixpoint of that function. That is if \(a\) is a pre-fixpoint
  of \(f\) then \(\funfont{lfp} \ f \leq a\).
\end{theorem}

\begin{theorem}\label{thm:lfp_le}
  If a point is less than or equal to any pre-fixpoint of a function
  then it is less than or equal to the least pre-fixpoint of that function.
  If for all pre-fixpoints \(b\) of \(f\) we have \(a \leq b\) then
  \(a \leq \funfont{lfp} \ f\).
\end{theorem}

We can now state the \emph{Knaster-Tarski} fixpoint theorem, which states that
the least fixpoint of a monotonic function on a complete lattice exists, and
is equal to the least pre-fixpoint of that function.

\begin{theorem}[Knaster-Tarski]\label{thm:lfp_eq}
  If \(f: \alpha \to \alpha\) is a monotonic function on a complete lattice
  then the least fixpoint of \(f\) exists and is equal to the least
  pre-fixpoint of \(f\). In other words:
  \[
    \funfont{lfp} \ f = f \ (\funfont{lfp} \ f)
  \]
\end{theorem}

\subsection{Set relations}

As we said before we will use the set theoretic view of relations, that is
a relation is a set of pairs.
We can, for example, define the \emph{identity relation}, as the the graph of the identity
function.

\begin{definition}
  The identity function graph \(\funfont{SRel.id}: \alpha \to_s \alpha\) is the 
  set of pairs \(\funfont{SRel.id} = \{(x, y) \mid x = y\}\).
\end{definition}

We can also define the composition of relations if there is a common element
between the codomain of the first relation and the domain of the second.

\begin{definition}
  The composition of two relations \(f: \alpha \to_s \beta\) and \(g: \beta \to_s \gamma\)
  is the set of pairs:
  \[
    f \circ g = \{(x, z) \mid \exists y, (x, y) \in f \land (y, z) \in g\}
  \]
\end{definition}

Of course we can prove that composing relations with the identity relation
is the same as the relation itself.

\begin{theorem}
  The composition of a relation with the identity relation is the relation itself.
  That is \(f \circ \funfont{id} = f\).
\end{theorem}

The same holds from the left. More importantly we can prove that the composition
is monotonic.

\begin{theorem}
  The composition of two functions is monotonic, that is if \(f \subseteq h\) and
  \(g \subseteq k\) then \(f \circ g \subseteq h \circ k\).
\end{theorem}

The last piece before defining the denotational semantics is the definition of
the \emph{set if then else}.

\begin{definition}
  The \emph{set if then else} (\(\funfont{Set.ite}\)) is a function that given three sets \(P, A, B\)
  returns the set:
  \[
    \funfont{Set.ite} \ P \ A \ B = A \cap P \cup B \setminus P
  \]
  where \(P\) is the \emph{predicate} or \emph{condition} set.
\end{definition}

\begin{theorem}
  The \(\funfont{Set.ite}\) function is monotonic in all its arguments.
  That is if \(a \subseteq b\) and \(a' \subseteq b'\) then
  \[
  \funfont{Set.ite} \ p \ a \ b \subseteq \funfont{Set.ite} \ p \ a' \ b'
  \]
\end{theorem}

\subsection{Denotational semantics}

We can finally define our function that given a command returns the 
\enquote{function} from state to state that represents the evaluation of that
command.

\begin{definition}[Denote]
We define the \(\funfont{denote}: \tyfont{Com} \to (\tyfont{State} \to_s \tyfont{State})\) 
function as:
  \begin{align*}
    \funfont{denote} \ \ImpSkip &\coloneqq \funfont{SRel.id} \\
    \funfont{denote} \ (v = a) &\coloneqq \{(s, t) \mid t = s[v \leftarrow a \ s]\} \\
    \funfont{denote} \ \ImpConcat{c}{c'} &\coloneqq \funfont{denote} \ c \circ \funfont{denote} \ c' \\
    \funfont{denote} \ (\ImpIfElse{b}{c}{c'}) &\coloneqq \funfont{Set.ite} \ \{(s, \placeholder) \mid b \ s\} \ (\funfont{denote} \ c) \ (\funfont{denote} \ c') \\
    \funfont{denote} \ (\ImpWhile{b}{c}) &\coloneqq \funfont{lfp} \ (W \ b \ (\funfont{denote} \ c))
  \end{align*}
  where \(W \ b \ f:\tyfont{Set} \ (\tyfont{State} \times \tyfont{State}) \to_o 
    \tyfont{Set} \ (\tyfont{State} \times \tyfont{State})\) is:
  \[
    W \ b \ f = \lambda g \mapsto \funfont{Set.ite} \ \{(s, \placeholder) \mid b \ s\} \ (f \circ g) \ \funfont{id}
  \]
\end{definition}

Of course this is only well defined if the least fixpoint of \(W \ b \ f\) exists,
for this it suffices to show that \(W \ b \ f\) is an order homomorphism, as we 
can then apply the \emph{Knaster-Tarski} theorem.

\begin{theorem}
  The function \(W \ b \ f\) is an order homomorphism.
\end{theorem}

Then, the fixpoint of \(W \ b \ f\)
is a \emph{relation} from the first state in which the \enquote{while} is
evaluated to the last state, that is, if the loop terminates. If the loop
doesn't terminate the relation is the \emph{empty set}, that is no state
is related to any other state.
Let's see a concrete example of what happens to the infinite loop 
\(L \coloneqq \ImpWhile{\texttt{true}}{\ImpSkip}\).
Applying the definition of the denotational semantics we have:
\[
  \funfont{denote} \ L = \funfont{lfp} \ (W \ \texttt{true} \ \funfont{SRel.id})
\]
unfolding the definition of \(W\) we have:
\begin{align*}
  W &= \lambda g \mapsto \funfont{SRel.id} \circ g \cap \funfont{Univ} \cup \funfont{SRel.id} \setminus \funfont{Univ} \\
    &= \lambda g \mapsto g = \funfont{id}
\end{align*}
Now, the least fixpoint of \(W\) is
\[
  \funfont{lfp} \ W = \bigcap \{a \mid \funfont{id} \ a \leq a\} = 
    \bigcap \funfont{Univ} = \emptyset
\]

\subsection{Kleene fixpoint theorem}

Defining the semantics with a relation works, but it makes computation difficult,
as we don't have a general way of computing the least fixpoint. We now 
introduce a theory, originally developed by \emph{Dana Scott}, that allows us
use \emph{partial functions} instead of relation. The problem with
partial functions is that they \emph{don't} form a complete lattice, as they
don't have an infimum. To solve this problem we will use the \emph{Kleene fixpoint}
theorem, which operates on \emph{\(\omega\)-continuous} functions.

We start by defining the concept of \(\omega\)-chain, which is a sequence of
sets where each set is a subset of the next.

\begin{definition}[\(\omega\)-chain]
  An \(\omega\)-chain is a sequence of sets \(\{S_i\}_{i \in \mathbb{N}}\) such that
  \(S_i \subseteq S_{i+1}\) for all \(i \in \mathbb{N}\).
\end{definition}

An \(\omega\)-continuous function is a function that \enquote{preserves the limit} of
\(\omega\)-chains.

\begin{definition}[\(\omega\)-continuous]
  A function \(f: \alpha \to \alpha\) is \(\omega\)-continuous if it commutes with the
  indexed union for every chain, that is for every \(\omega\)-chain 
  \(\{S_i\}_{i \in \mathbb{N}}\) we have:
  \[
    f \left( \bigcup_{i \in \mathbb{N}} S_i \right) = \bigcup_{i \in \mathbb{N}} f \ S_i
  \]
\end{definition}

\begin{notation}
  We will denote an \(\omega\)-continuous function as \(f: \alpha \to_\omega \alpha\).
\end{notation}

Continuity is a \enquote{stronger} property than monotonicity, as every continuous
function is monotonic, but not every monotonic function is continuous.

\begin{theorem}
  Every \(\omega\)-continuous function is monotonic.
\end{theorem}

As an example of a \(\omega\)-chain we consider the sequence \emph{power function} of
an order homomorphism from sets to sets.

\begin{theorem}\label{thm:fpow_chain}
  The sequence of power functions of an order homomorphism \(f: \tyfont{Set} \ \alpha \to_o \tyfont{Set} \ \alpha\)
   is an \(\omega\)-chain. That is \(\{f^n \ \emptyset\}_{n \in \mathbb{N}}\) is an \(\omega\)-chain.
\end{theorem}

This will be useful later, as we prove that the least fixpoint of an \(\omega\)-continuous
function is a \emph{fixed point}.

\begin{definition}[Least fixpoint]
  The least fixpoint of a function \(f: \alpha \to_\omega \alpha\) is the 
  \enquote{limit} of the sequence of the power functions of \(f\), that is:
  \[
    \funfont{lfp} \ f = \bigcup_{n \in \mathbb{N}} f^n \ \emptyset
  \]
\end{definition}

Now we can state the \emph{Kleene fixpoint theorem}.

\begin{theorem}[Kleene fixpoint]
  If \(f: \alpha \to_\omega \alpha\) is an \(\omega\)-continuous function
  then the least fixpoint of \(f\) exists and is equal to the least fixpoint
  of \(f\) as an order homomorphism.
\end{theorem}

\begin{proof}
  We have to prove that \(\funfont{OrderHom.lfp f = ContinuousHom.lfp f}\),
  we do this by proving that both sets contain the other.
  \begin{itemize}
    \item[(\(\subseteq\))] Suffices to show that \(\bigcup_{n \in \mathbb{N}} f^n \ \emptyset\)
      is a pre-fixpoint of \(f\), as we can then apply~\ref{thm:lfp_le} obtaining
      the result. That the leastt fixpoint is a pre-fixpoint is more or less
      immediate by using~\ref{thm:fpow_chain}.

    \item[(\(\supseteq\))] For this it suffices: \(\funfont{lfp} \ f \supseteq f^n \ \emptyset\),
      we proceed by induction on \(n\). The base case is trivially contradictory,
      the inductive step is immediate by applying Knaster-Tarski~\ref{thm:lfp_eq}
      as \(f\) is an order homomorphism.
  \end{itemize}
\end{proof}

This fixpoint definition has the advantage of giving us an easy approximation
of the least fixpoint, as we can take the union of the first \(n\) power functions
of \(f\).

\section{Equivalence of semantics}

% Appendices 

% appendix
\appendix

\section{Normalization in Lean}
\subsection{Omega combinator}

\section{Equality: rec, ndrec and subst}

\section{\texttt{IMP} Turing completeness}

\section{Future work}

Things to add:

\begin{itemize}
  \item Add formalization of \(\lambda\)-calculus and proofs.
  \item Compiler.
  \item Typed \texttt{IMP}.
  \item Hoare Logic.
  \item Abstract interpretation.
\end{itemize}

\section{Código}

El código del texto se encuentra en el repositorio de GitHub
\url{https://github.com/remind-me-later/SemanticsLean}. 
Lo reproducimos aquí por completitud, con una subsección para cada
archivo.

\subsection{Mapas}

\inputminted{lean4}{../SemanticsLean/Semantics/Maps.lean}

\subsection{Conjuntos}

\inputminted{lean4}{../SemanticsLean/Semantics/Set.lean}

\subsection{Sintaxis de IMP}

\inputminted{lean4}{../SemanticsLean/Semantics/Imp/Lang.lean}

\subsection{Expresiones aritméticas}

\inputminted{lean4}{../SemanticsLean/Semantics/Imp/Aexp.lean}

\subsection{Expresiones booleanas}

\inputminted{lean4}{../SemanticsLean/Semantics/Imp/Bexp.lean}

\subsection{Semántica de paso largo}

\inputminted{lean4}{../SemanticsLean/Semantics/Imp/BigStep.lean}

\subsection{Clausura reflexiva transitiva}

\inputminted{lean4}{../SemanticsLean/Semantics/ReflTransRel.lean}

\subsection{Semántica operacional}

\inputminted{lean4}{../SemanticsLean/Semantics/Imp/SmallStep.lean}

\subsection{Orden y punto fijo de Knaster-Tarski}

\inputminted{lean4}{../SemanticsLean/Semantics/Lattice.lean}

\subsection{Cadenas y punto fijo de Kleene}

\inputminted{lean4}{../SemanticsLean/Semantics/Chain.lean}

\subsection{Semántica denotacional}

\inputminted{lean4}{../SemanticsLean/Semantics/Imp/Denotational.lean}

\subsection{Equivalencia de semánticas}

\inputminted{lean4}{../SemanticsLean/Semantics/Imp/Equivalence.lean}

% Bibliography

\bibliographystyle{alpha}
\bibliography{refs}


\end{document}
